\chapter{Current Mirrors}\label{ch:Current_Mirrors}

This chapter covers various types of current mirrors,
which are circuits designed to copy (or ``mirror'') a current
from one of its branches to another (amplifying or attenuating
it in the process).

Ideally a current mirror should have perfect current
replication, infinite output resistance, zero input
resistance and no voltage drop across the mirror. However,
real current mirrors deviate from the ideal due to various
non-ideal effects, including component tolerances, component
parasitics, loading effects, channel length modulation,
base width modulation and temperature variations.

Channel Length Modulation and Base Width Modulation are
worth a discussion.
\begin{enumerate}
  \item \textbf{(MOSFET) Channel Length Modulation:} Generally
    we assume $K_{CLM} = 0$ in MOSFET equations when analyzing
    transistor behaviour for simplicity. In real circuits,
    however, this is not the case. When $K_{CLM}$ is non-zero,
    the drain current varies with changes in the drain-source
    voltage by a factor of $(1 + K_{CLM} V_{DS})$. This affects
    the accuracy of current mirroring as the output current
    will vary with changes in the drain to source voltage.
    Cascoding may help reduce the impact of channel length modulation.
  \item \textbf{(BJT) Base Width Modulation:} Similar to channel length
    modulation in MOSFETs, BJTs also exhibit a phenomenon
    known as base width modulation (or Early effect). This
    effect causes the collector current to vary with changes
    in the collector-emitter voltage, leading to inaccuracies
    in current mirroring.
\end{enumerate}

All mirrors in this chapter are implemented using N-Channel
MOSFETs. It is possible to adapt the configurations to
BJTs, or other transistor types regardless of technology.
Refer to~\ref{sec:theory_general_complementary} for more
information on converting circuits to their complementary
forms.

Every mirror is analysed in two passes (Chapter~\ref{ch:analysis}):
Large Signal (Q5) with the square law of
Section~\ref{sec:devices_mosfet} to set bias currents and
voltages, then Small Signal (Q4) with \(g_m\), \(g_o\) to find
current gain, \(Z_{in}\) and \(Z_{out}\). Resistors that set the
input current are kept in the small signal circuit (\(G\) terms),
since they load the input node.

The primary use of current mirrors is to provide a stable
current source or sink in analog circuits. They can also
be used to steer currents into different branches of a
circuit, or to replace resistive loads in amplifiers and
other circuits thereby acting as active loads.

\section{Topologies}
\subsection{Basic Mirror}\label{sec:mirror_basic}

\subsubsection{Overview}
Constructed from two transistors, a diode connected
\(M_0\) and a mirroring \(M_1\) sharing \(V_{GS}\). It is the
simplest form of a current mirror; its current gain is set
by the ratio of \(\frac{W}{L}\). (This block was previously
mislabelled as the Widlar mirror, which adds a source
resistor, see Section~\ref{sec:mirror_widlar}.)

\subsubsection{Schematic}
\begin{circuitfig}
  \draw (2.98, 7.905) to[american voltage source, l_={$V_{cc}$}] (1.73, 7.905);
	\node[nmos, arrowmos, xscale=-1] (N1) at (3.25, 4.405){} node[anchor=east] at (N1.text){$M_0$};
	\node[nmos, arrowmos] (N2) at (5.75, 4.405){} node[anchor=west] at (N2.text){$M_1$};
	\draw (3.25, 7.405) to[european resistor, l={$G_2$}] (3.25, 5.905);
	\draw (5.75, 5.655) -| (5.75, 5.175);
	\draw (4.5, 4.405) |- (3.25, 5.405);
	\draw (4.77, 4.405) -- (4.23, 4.405);
	\draw (3.25, 3.635) -| (3.25, 3.405) -| (5.75, 3.635);
	\draw (3.75, 2.655) to[american voltage source, l={$V_{ee}$}] (2.5, 2.655);
	\draw (3.75, 2.655) -| (4.5, 3.405);
	\node[ground, rotate=-90] at (1.73, 7.905){};
	\node[ground, rotate=-90] at (2.5, 2.655){};
	\draw[dash pattern={on 1.6pt off 0.8pt}, -stealth] (5.75, 5.655) -- (6.25, 5.655);
	\node[shape=rectangle, minimum width=0.965cm, minimum height=0.465cm] (N3) at (6.75, 5.615){} node[anchor=center] at (N3.text){$V_{out}$};
	\draw (5.75, 7.365) to[european resistor, l={\textcolor{rgb,255:red,255;green,0;blue,0}{$G_3$}}] (5.75, 5.865);
	\draw (5.75, 5.655) -- (5.75, 5.885);
	\draw (3.25, 7.405) |- (2.98, 7.905);
	\draw (3.25, 7.905) -| (5.75, 7.365);
	\node[shape=rectangle, minimum width=4.965cm, minimum height=1.09cm] at (4.375, 9.928){} node[anchor=north, align=center, inner sep=6pt] at (4.375, 10.49){Large Signal Schematic\\(Sink)};
	\node[shape=rectangle, draw={rgb,255:red,255;green,0;blue,0}, line width=1pt, dash pattern={on 4pt off 2pt on 1pt off 2pt on 1pt off 2pt}, minimum width=1.465cm, minimum height=2.195cm] at (6, 7){};
	\node[shape=rectangle, minimum width=3.09cm, minimum height=2.465cm] at (8.438, 6.865){} node[anchor=north west, align=left, inner sep=6pt] at (6.875, 8.115){\textcolor{rgb,255:red,255;green,0;blue,0}{\footnotesize Exists Only For Analysis. In  Implementation take output directly from V\_\{out\}}};
	\node[shape=rectangle, minimum width=0.965cm, minimum height=0.465cm] (N4) at (2.25, 5.405){} node[anchor=center] at (N4.text){$V_{x}$};
	\draw[dash pattern={on 1.6pt off 0.8pt}, -stealth] (3.25, 5.405) |- (2.75, 5.405);
	\node[shape=rectangle, minimum width=0.965cm, minimum height=0.465cm] (N5) at (4.5, 6.115){} node[anchor=center] at (N5.text){$I_{G}$};
	\draw (3.25, 5.175) -| (3.25, 5.405);
	\draw (3.25, 5.405) -| (3.25, 5.905);
	\draw (10, 1.75) |- (10, 11);
	\draw (17.25, 5.563) to[american current source, l_={$g_{m_1}V_{gs}$}] (17.25, 4.063);
	\draw (18.75, 5.563) to[european resistor, l_={$g_{o_1}$}] (18.75, 4.063);
	\draw (17.25, 4.063) -- (18.75, 4.063);
	\draw (18.75, 5.563) -- (17.25, 5.563);
	\draw (17.25, 4.063) -- (16, 4.063);
	\draw (16, 5.563) -- (15.25, 5.563);
	\node[shape=rectangle, minimum width=0.465cm, minimum height=0.465cm] at (16, 5.313){} node[anchor=center, align=center, inner sep=6pt] at (16, 5.313){\footnotesize G};
	\node[shape=rectangle, minimum width=0.465cm, minimum height=0.465cm] at (16, 4.313){} node[anchor=center, align=center, inner sep=6pt] at (16, 4.313){\footnotesize S};
	\node[shape=rectangle, minimum width=0.465cm, minimum height=0.465cm] at (19.25, 5.313){} node[anchor=center, align=center, inner sep=6pt] at (19.25, 5.313){\footnotesize D};
	\draw (15.25, 5.563) -- (14.5, 5.563);
	\draw (13, 5.563) to[american current source, l={$g_{m_0}V_{gs}$}] (13, 4.063);
	\draw (11.5, 5.563) to[european resistor, l={$g_{o_0}$}] (11.5, 4.063);
	\draw (13, 5.563) -- (11.5, 5.563);
	\draw (13, 4.063) -- (11.5, 4.063);
	\draw (14.5, 4.063) -- (13, 4.063);
	\node[shape=rectangle, minimum width=0.465cm, minimum height=0.465cm] at (14.5, 5.313){} node[anchor=center, align=center, inner sep=6pt] at (14.5, 5.313){\footnotesize G};
	\node[shape=rectangle, minimum width=0.465cm, minimum height=0.465cm] at (14.5, 4.313){} node[anchor=center, align=center, inner sep=6pt] at (14.5, 4.313){\footnotesize S};
	\node[shape=rectangle, minimum width=0.465cm, minimum height=0.465cm] at (11, 5.313){} node[anchor=center, align=center, inner sep=6pt] at (11, 5.313){\footnotesize D};
	\draw (11.5, 7.563) to[european resistor, l={$G_2$}] (11.5, 6.063);
	\draw (11.5, 6.063) -| (11.5, 5.562);
	\draw (11.5, 7.563) -| (11.5, 8.062) -- (12.25, 8.062) -| (12.25, 7.812);
	\draw (18.75, 7.563) to[european resistor, l={\textcolor{rgb,255:red,255;green,0;blue,0}{$G_3$}}] (18.75, 6.063);
	\draw (18.75, 6.063) -| (18.75, 5.562);
	\draw (18.75, 7.563) -| (18.75, 8.062) |- (18, 8.062) -| (18, 7.812);
	\draw[dash pattern={on 1.6pt off 0.8pt}, -stealth] (18.75, 5.563) -- (19.25, 5.563);
	\node[shape=rectangle, minimum width=0.715cm, minimum height=0.465cm] (N6) at (10.625, 5.563){} node[anchor=center] at (N6.text){$V_{x}$};
	\node[shape=rectangle, minimum width=0.965cm, minimum height=0.465cm] (N7) at (19.75, 5.563){} node[anchor=center] at (N7.text){$V_{out}$};
	\node[ground] at (14.5, 4.063){};
	\node[ground] at (16, 4.063){};
	\node[ground] at (12.25, 7.812){};
	\node[ground] at (18, 7.812){};
	\node[shape=rectangle, minimum width=0.965cm, minimum height=0.465cm] (N8) at (15.25, 6.813){} node[anchor=center] at (N8.text){$I_{G}$};
	\draw[dash pattern={on 1.6pt off 0.8pt}, -stealth] (11.5, 5.563) -- (11, 5.563);
	\node[shape=rectangle, minimum width=4.215cm, minimum height=1.09cm] at (15.375, 9.937){} node[anchor=north, align=center, inner sep=6pt] at (15.375, 10.5){Small Signal Schematic\\(Sink)};
	\node[circ] at (3.25, 5.405){};
	\node[circ] at (4.5, 4.405){};
	\node[circ] at (4.5, 3.405){};
	\node[circ] at (3.25, 7.905){};
	\draw (13, 5.563) |- (15.25, 6.063) -| (15.25, 5.563);
	\draw[dash pattern={on 1.6pt off 0.8pt}, -stealth] (15.25, 6.063) -- (15.25, 6.563);
	\draw[dash pattern={on 1.6pt off 0.8pt}, -stealth] (4.5, 5.365) -- (4.5, 5.865);
	\node[circ] at (11.5, 5.563){};
	\node[circ] at (15.25, 5.563){};
	\node[circ] at (18.75, 5.563){};
	\node[circ] at (13, 5.563){};
	\node[shape=rectangle, draw={rgb,255:red,255;green,0;blue,0}, line width=1pt, dash pattern={on 4pt off 2pt on 1pt off 2pt on 1pt off 2pt}, minimum width=1.465cm, minimum height=1.965cm] at (19, 7.063){};
\end{circuitfig}

\subsubsection{Transfer Function}
Here's a table of the voltages at each terminal of the transistors:

\noindent\begin{tabularx}{\linewidth}{|X|X|X|X|X|X|} \hline
  Transistor & Gate (G) & Source (S) & Drain (D) & $V_{GS}$ & $V_{DS}$ \\ \hline
  $M_0$ & $V_{X}$ & $V_{ee}$ & $V_{X}$ & $V_{X} - V_{ee}$ & $V_{X} - V_{ee}$ \\ \hline
  $M_1$ & $V_{X}$ & $V_{ee}$ & $V_{OUT}$ & $V_{X} - V_{ee}$ & $V_{OUT} - V_{ee}$ \\ \hline
\end{tabularx}

\begin{enumerate}
  \item \textbf{Large Signal, Time Domain, DC Input Analysis (Q5)}:\\
    Applying KCL at $V_{OUT}$:
    \begin{align*}
      (V_{OUT} - V_{cc})G_3 + I_{M_1} &= 0 \\
      V_{OUT}G_3 - V_{cc}G_3 + I_{M_1} &= 0 \\
      V_{OUT} &= V_{cc} - \frac{I_{M_1}}{G_3} \tag{1}\end{align*}
    Applying KCL at $V_{X}$:
    \begin{align*}
      (V_{X} - V_{cc})G_2 + I_{M_0} + I_{G} &= 0 \\
      V_{X}G_2 - V_{cc}G_2 + (I_{M_0} + I_{G}) &= 0 \\
      V_{X}G_2 &= - (I_{M_0} + I_{G}) + V_{cc}G_2 \\
      V_{X} &= V_{cc} - \frac{(I_{M_0} + I_{G})}{G_2} \tag{2}
    \end{align*}
    Using the transistor equations for $I_{M_0}$ and
    $I_{M_1}$:
    \begin{align*}
      I_{M_0} &= k_{n_0}(V_{GS_0} - V_{th})^2 \cdot (1 + K_{CLM} V_{DS_0}) \\
              &= k_{n_0}(V_{X} - V_{ee} - V_{th})^2 \cdot (1 + K_{CLM} (V_{X} - V_{ee})) \\
              &= k_{n_0}\left(\left(V_{cc} - \frac{(I_{M_0} + I_{G})}{G_2}\right) - V_{ee} - V_{th}\right)^2 \cdot \left(1 + K_{CLM} \left(\left(V_{cc} - \frac{(I_{M_0} + I_{G})}{G_2}\right) - V_{ee}\right)\right) \\
              &= k_{n_0}\left(\left(V_{cc} - \frac{(I_{M_0} + I_{G})}{G_2}\right) - V_{ee} - V_{th}\right)^2 \quad \text{[assuming } K_{CLM} = 0 \text{]} \tag{3} \\ \\
      I_{M_1} &= k_{n_1}(V_{GS_1} - V_{th})^2 \cdot (1 + K_{CLM} V_{DS_1}) \\
              &= k_{n_1}(V_{X} - V_{ee} - V_{th})^2 \cdot (1 + K_{CLM} (V_{OUT} - V_{ee})) \\
              &= k_{n_1}\left(\left(V_{cc} - \frac{(I_{M_0} + I_{G})}{G_2}\right) - V_{ee} - V_{th}\right)^2 \cdot \left(1 + K_{CLM} \left(V_{cc} - \frac{I_{M_1}}{G_3} - V_{ee}\right)\right) \\
              &= k_{n_1}\left(\left(V_{cc} - \frac{(I_{M_0} + I_{G})}{G_2}\right) - V_{ee} - V_{th}\right)^2 \quad \text{[assuming } K_{CLM} = 0 \text{]} \tag{4}
    \end{align*}
    The input current is the current through \(G_2\) and the output
    current is the drain current of \(M_1\):
    \begin{align*}
      I_{IN} &= (V_{cc} - V_{X})G_2 = I_{M_0} + I_{G} \approx I_{M_0} \quad \text{[from (2), } I_G \approx 0 \text{ by MOS physics]} \\
      I_{OUT} &= (V_{cc} - V_{OUT})G_3 = I_{M_1} \quad \text{[from (1)]}
    \end{align*}
    Both transistors share \(V_{GS} = V_X - V_{ee}\), so dividing the
    second line of (4) by the second line of (3) cancels the square term:
    \begin{align*}
      \frac{I_{M_1}}{I_{M_0}} &= \frac{k_{n_1}(V_{X} - V_{ee} - V_{th})^2 (1 + K_{CLM} (V_{OUT} - V_{ee}))}{k_{n_0}(V_{X} - V_{ee} - V_{th})^2 (1 + K_{CLM} (V_{X} - V_{ee}))} \\
                              &= \frac{k_{n_1}}{k_{n_0}} \cdot \frac{1 + K_{CLM} (V_{OUT} - V_{ee})}{1 + K_{CLM} (V_{X} - V_{ee})} \\
                              &= \frac{k_{n_1}}{k_{n_0}} \quad \text{[assuming } K_{CLM} = 0 \text{, or } V_{OUT} = V_X \text{]}
    \end{align*}
    Thus, with \(k_n = \frac{K_P}{2}\frac{W}{L}\),
    \begin{align*}
      \frac{I_{OUT}}{I_{IN}} \approx \frac{I_{M_1}}{I_{M_0}} \approx \frac{k_{n_1}}{k_{n_0}} = \frac{\frac{K_P}{2} \frac{W_1}{L_1}}{\frac{K_P}{2} \frac{W_0}{L_0}} = \frac{\frac{W_1}{L_1}}{\frac{W_0}{L_0}} = \frac{W_1 L_0}{L_1 W_0}
    \end{align*}
  \item \textbf{Small Signal, Frequency Domain, AC Input Analysis (Q4)}:\\
    Applying KCL at $V_{out}$:
    \begin{align*}
      V_{out}G_3 + V_{out}g_{o_1} + g_{m_1}V_{x} &= 0 \\
      V_{out}(G_3 + g_{o_1}) + g_{m_1}V_{x} &= 0 \tag{5}
    \end{align*}
    Applying KCL at $V_{x}$:
    \begin{align*}
      V_{x}G_2 + V_{x}g_{o_0} + g_{m_0}V_{x} + I_{G} &= 0 \\
      V_{x}(G_2 + g_{o_0} + g_{m_0}) &= 0 \quad \text{[assuming } I_G \approx 0 \text{ by MOS physics]} \tag{6}
    \end{align*}
    Adding voltage source $V_s$ at $V_x$ that supplies
    a current $I_s$ into $V_{x}$, we get:
    \begin{align*}
      V_{x}(G_2 + g_{o_0} + g_{m_0}) - I_{s} &=0 \tag{7} \\
      V_{x} &= V_{s} \tag{8}
    \end{align*}
    Solving equations (5), (7) and (8) ((7) only fixes \(I_s\)):
    \begin{align*}
      V_{out}(G_3 + g_{o_1}) &= - g_{m_1} V_{x} \quad \text{[from (5)]} \\
      V_{out} &= - \frac{g_{m_1} V_{x}}{G_3 + g_{o_1}} \\
              &= - \frac{g_{m_1} V_{s}}{G_3 + g_{o_1}} \quad \text{[from (8)]} \\
      H(s) = \frac{N(s)}{D(s)} = \frac{V_{out}}{V_{s}} &= - \frac{g_{m_1}}{G_3 + g_{o_1}}
    \end{align*}
\end{enumerate}

\subsubsection{Analysis}
\begin{enumerate}
  \item \textbf{Channel Length Modulation, Large Signal, Time Domain, DC Input (Q5)}:
    the ratio \(\frac{W_1 L_0}{L_1 W_0}\) above used \(K_{CLM} = 0\), but \(Z_{out}\) below
    is \(\frac{1}{g_{o_1}}\) with \(g_{o_1} = \frac{K_{CLM} I_{M_1}}{1 + K_{CLM} V_{DS_1}} \neq 0\). Both come from the same
    \(K_{CLM}\), so it is kept here to first order. From the exact ratio:
    \begin{align*}
      \frac{I_{M_1}}{I_{M_0}} &= \frac{k_{n_1}}{k_{n_0}} \cdot \frac{1 + K_{CLM} (V_{OUT} - V_{ee})}{1 + K_{CLM} (V_{X} - V_{ee})} \\
      &= \frac{k_{n_1}}{k_{n_0}} \cdot \frac{1 + K_{CLM} (V_{X} - V_{ee}) + K_{CLM} (V_{OUT} - V_{X})}{1 + K_{CLM} (V_{X} - V_{ee})} \hspace{1cm} [\pm K_{CLM} V_X] \\
      &= \frac{k_{n_1}}{k_{n_0}} \left(1 + \frac{K_{CLM} (V_{OUT} - V_{X})}{1 + K_{CLM} (V_{X} - V_{ee})}\right)
    \end{align*}
    The ratio is exact only at \(V_{OUT} = V_X\) (\(V_{DS_1} = V_{DS_0}\)). Differentiating
    \(I_{M_1}\) with respect to \(V_{OUT}\) at fixed \(I_{M_0}\):
    \begin{align*}
      \frac{dI_{M_1}}{dV_{OUT}} &= \frac{k_{n_1}}{k_{n_0}} I_{M_0} \cdot \frac{K_{CLM}}{1 + K_{CLM} (V_{X} - V_{ee})}
      = \frac{K_{CLM} I_{M_1}|_{V_{OUT} = V_X}}{1 + K_{CLM} (V_{X} - V_{ee})} = g_{o_1}|_{V_{OUT} = V_X}
    \end{align*}
    so the Q5 error and the Q4 output conductance are the same effect:
    \(I_{OUT} \approx \frac{W_1 L_0}{L_1 W_0} I_{IN} + g_{o_1}(V_{OUT} - V_X)\). Setting \(K_{CLM} = 0\) is the
    leading order of the ratio; \(g_{o_1}\) is the leading order of its slope.
  \item \textbf{Output Load (Max), Large Signal, Time Domain, DC Input (Q5)}:
    \(M_1\) stays saturated while \(V_{DS_1} \geq V_{ov_1}\). Both
    transistors share \(V_{GS}\) and \(V_{th}\), so
    \(V_{ov_1} = V_{GS_1} - V_{th} = V_{GS_0} - V_{th} = V_{ov_0}\).
    With the load \(R_L = \frac{1}{G_3}\) to \(V_{cc}\), (1) gives
    \(V_{OUT} = V_{cc} - I_{M_1} R_L\):
    \begin{align*}
      V_{DS_1} = V_{OUT} - V_{ee} &\geq V_{ov_1} \\
      V_{OUT,\min} &= V_{ee} + V_{ov_1} \\
      V_{cc} - I_{M_1} R_L &\geq V_{OUT,\min} \\
      I_{M_1} R_L &\leq V_{cc} - V_{OUT,\min} \\
      R_{L, \max} &= \frac{V_{cc} - V_{OUT,\min}}{I_{M_1}} = \frac{V_{L,\max}}{I_{M_1}} \quad \text{[} V_{L,\max} = V_{cc} - V_{OUT,\min} \text{]}
    \end{align*}
  \item \textbf{Current Gain, Small Signal, Frequency Domain, AC Input (Q4)}:
    Inject \(I_{in}\) into \(V_x\) and short the output to AC ground
    (\(V_{out} = 0\)). (6) with the injected current, and the drain
    current of \(M_1\) flowing from the AC ground into \(V_{out}\):
    \begin{align*}
      V_{x}(G_2 + g_{o_0} + g_{m_0}) &= I_{in} \\
      V_{x} &= \frac{I_{in}}{G_2 + g_{o_0} + g_{m_0}} \\
      I_{out} &= g_{m_1} V_x + g_{o_1} V_{out} = g_{m_1} V_x + g_{o_1} \cdot 0 = g_{m_1} V_x \\
              &= \frac{g_{m_1}}{G_2 + g_{o_0} + g_{m_0}} I_{in} \\
      \frac{I_{out}}{I_{in}} &= \frac{g_{m_1}}{G_2 + g_{o_0} + g_{m_0}} \approx \frac{g_{m_1}}{g_{m_0}} \quad \text{[} G_2, g_{o_0} \ll g_{m_0} \text{]}
    \end{align*}
    With \(g_m \approx \sqrt{2 K_P \frac{W}{L} I_D}\)
    (Section~\ref{sec:devices_mosfet}) and the Q5 result
    \(I_{M_1} = \frac{W_1 L_0}{L_1 W_0} I_{M_0}\):
    \begin{align*}
      \frac{g_{m_1}}{g_{m_0}} &= \frac{\sqrt{2K_P \frac{W_1}{L_1} I_{M_1}}}{\sqrt{2K_P \frac{W_0}{L_0} I_{M_0}}}
                               = \sqrt{\frac{\frac{W_1}{L_1}}{\frac{W_0}{L_0}} \cdot \frac{I_{M_1}}{I_{M_0}}} \\
                              &= \sqrt{\frac{W_1 L_0}{L_1 W_0} \cdot \frac{W_1 L_0}{L_1 W_0}} = \frac{W_1 L_0}{L_1 W_0}
    \end{align*}
  \item \textbf{Impedances}:
    \paragraph{At Input:} Applying a test current $I_{test}$
      at $V_{x}$, from equation (6):
      \begin{align*}
        V_{x}(G_2 + g_{o_0} + g_{m_0}) &= I_{test} \\
        \therefore Z_{in} &= \frac{V_{x}}{I_{test}} = \frac{1}{G_2 + g_{o_0} + g_{m_0}}
      \end{align*}
      Ignoring $G_2$ and assuming $g_{m_0} \gg g_{o_0}$, with
      \(g_{m_0} = \frac{2 I_{M_0}}{V_{ov_0}}\), we get:
      \begin{align*}
        Z_{in} &\approx \frac{1}{g_{m_0}} = \frac{1}{\frac{2 I_{M_0}}{V_{ov_0}}} = \frac{V_{ov_0}}{2 I_{M_0}}
      \end{align*}
    \paragraph{At Output:} Applying a test current $I_{test}$
      at $V_{out}$, from equation (5):
      \begin{align*}
        V_{out}(G_3 + g_{o_1}) + g_{m_1}V_{x} &= I_{test} \\
        V_{out}(G_3 + g_{o_1}) + g_{m_1}(0) &= I_{test} \quad \text{[} V_x = 0 \text{ from (6)]} \\
        \therefore Z_{out} &= \frac{V_{out}}{I_{test}} = \frac{1}{G_3 + g_{o_1}}
      \end{align*}
      Ignoring $G_3$, with \(g_{o_1} \approx K_{CLM} I_{M_1}\), we get:
      \begin{align*}
        Z_{out} &\approx \frac{1}{g_{o_1}} = \frac{1}{K_{CLM} I_{M_1}}
      \end{align*}
\end{enumerate}

\subsubsection{Design}
Large Signal, Time Domain, DC Input Analysis (Q5) can be used to set the bias
currents and voltages in this circuit.
\begin{enumerate}
  \item Choose appropriate transistor dimensions
    ($W$ and $L$) for transistors $M_0$ and $M_1$
    to achieve the desired current gain.
    \begin{align*}
      \frac{I_{OUT}}{I_{IN}} &= \frac{W_1 L_0}{L_1 W_0}
    \end{align*}
  \item Choose a suitable overdrive voltage value
    ($V_{ov}$) for computing the bias voltage at
    node $V_{X}$.
    \begin{align*}
      V_{ov_0} &= V_{GS_0} - V_{th} \\
      V_{ov_0} &= (V_{X} - V_{ee}) - V_{th} \\
      V_{X} &= V_{ov_0} + V_{th} + V_{ee}
    \end{align*}
  \item Using the chosen $V_{ov}$, compute $W_0$
    to set the required bias current $I_{IN}$.
    \begin{align*}
      I_{IN} = I_{M_0} &= \frac{K_P}{2} \frac{W_0}{L_0} V_{ov_0}^2 \\
      \therefore W_0 &= \frac{2 I_{IN} L_0}{K_P V_{ov_0}^2}
    \end{align*}
    $W_1$ can then be computed using the current
    gain relation from step (a).
  \item Compute the required value of $G_2$ (or
    $R_2$) to ensure that the voltage at node
    $V_{X}$ is as calculated in step (b).
    \begin{align*}
      (V_{X} - V_{cc})G_2 + I_{M_0} + I_{G} &= 0 \\
      (V_{X} - V_{cc})G_2 + I_{M_0} &\approx 0 \quad \text{[assuming } I_G \approx 0 \text{ by MOS physics ]} \\
      (V_{X} - V_{cc})G_2 &\approx - I_{M_0} \\
      G_2 &\approx - \frac{I_{M_0}}{(V_{X} - V_{cc})} \\
      \therefore G_2 &= \frac{I_{M_0}}{V_{cc} - V_{X}} \\
      R_2 &= \frac{1}{G_2} = \frac{V_{cc} - V_{X}}{I_{M_0}}
    \end{align*}
\end{enumerate}

\subsubsection{Question}
Design a Basic Current Mirror (Sink) that meets
the following requirements. The mirror must provide
a current gain at the output of 5. The input current
must be $I_{IN} = 100 \mu A$. The mirror must operate
with $V_{cc} = 5V$ and $V_{ee} = 0V$. Assume the
following transistor parameters: $V_{th} = 1V$,
$K_{CLM} = 0.02 V^{-1}$, $K_P = 200 \mu A/V^2$ and
$L = 1 \mu m$ for all transistors. Find the input and
output impedance of the mirror and the maximum load
that can be connected at the output.
\paragraph{Solution:}
\begin{enumerate}
  \item \textbf{Choose appropriate transistor dimensions:}
    \begin{align*}
      \frac{I_{OUT}}{I_{IN}} &= \frac{W_1 L_0}{L_1 W_0} = 5 \\
      \therefore W_1 &= 5 W_0 \quad \text{[assuming } L_0 = L_1 \text{]}
    \end{align*}
  \item \textbf{Choose a suitable overdrive voltage and determine
    the resulting bias voltages:}
    Choosing \(V_{ov_0} = 0.2V\):
    \begin{align*}
      V_{X} &= V_{ov_0} + V_{th} + V_{ee} = 0.2V + 1V + 0V = 1.2V
    \end{align*}
  \item \textbf{Compute $W_0$ to set the required bias current:} For $I_{IN} = 100 \mu A$
    \begin{align*}
      W_0 &= \frac{2 I_{IN} L_0}{K_P V_{ov_0}^2} \\
          &= \frac{2 \cdot 100 \mu A \cdot 1 \mu m}{200 \mu A/V^2 \cdot (0.2V)^2} \\
          &= \frac{200 \times 10^{-12} \, A \cdot m}{8 \times 10^{-6} \, A} \\
          &= 25 \mu m \\
      \therefore W_1 &= 5 W_0 = 5 \cdot 25 \mu m = 125 \mu m
    \end{align*}
  \item \textbf{Compute the required value of $G_2$ (or $R_2$):}
    \begin{align*}
      R_2 &= \frac{1}{G_2} = \frac{V_{cc} - V_{X}}{I_{M_0}} \\
          &= \frac{5V - 1.2V}{100 \mu A} \\
          &= 38 k\Omega
    \end{align*}
  \item \textbf{Compute the input impedance:}
    \begin{align*}
      g_{m_0} &= \frac{2 I_{M_0}}{V_{ov_0}} = \frac{2 \cdot 100 \mu A}{0.2V} = 1 mS \hspace{1cm}
      g_{o_0} \approx K_{CLM} I_{M_0} = 0.02 V^{-1} \cdot 100 \mu A = 2 \mu S \\
      Z_{in} &\approx \frac{1}{g_{m_0}} = \frac{V_{ov_0}}{2 I_{M_0}} = \frac{0.2V}{2 \cdot 100 \mu A} = 1 k\Omega
    \end{align*}
    Keeping \(G_2 = \frac{1}{38 k\Omega} = 26.3 \mu S\) and \(g_{o_0}\):
    \begin{align*}
      Z_{in} &= \frac{1}{G_2 + g_{o_0} + g_{m_0}} = \frac{1}{26.3 \mu S + 2 \mu S + 1000 \mu S} = \frac{1}{1028.3 \mu S} = 972 \Omega
    \end{align*}
  \item \textbf{Compute the output impedance:} The output current is
    \(I_{M_1} = 5 I_{IN} = 5 \cdot 100 \mu A = 500 \mu A\).
    \begin{align*}
      Z_{out} &\approx \frac{1}{g_{o_1}} \approx \frac{1}{K_{CLM} I_{M_1}} \\
              &= \frac{1}{0.02 V^{-1} \cdot 500 \mu A} \\
              &= 100 k\Omega
    \end{align*}
  \item \textbf{Actual current gain} (\(K_{CLM} = 0.02\,V^{-1}\) kept,
    \(V_{DS_0} = V_X - V_{ee} = 1.2V\)):
    \begin{align*}
      \frac{I_{M_1}}{I_{M_0}} &= 5 \cdot \frac{1 + 0.02\,(V_{OUT} - 0)}{1 + 0.02 \cdot 1.2} = 5 \cdot \frac{1 + 0.02\,V_{OUT}}{1.024} \\
      V_{OUT} = 1.2V &: \hspace{0.5cm} 5 \cdot \frac{1.024}{1.024} = 5.00 \\
      V_{OUT} = 0.2V &: \hspace{0.5cm} 5 \cdot \frac{1.004}{1.024} = 4.90 \\
      V_{OUT} = 5V &: \hspace{0.5cm} 5 \cdot \frac{1.1}{1.024} = 5.37
    \end{align*}
    The slope is \(\frac{5 \cdot 100\,\mu A \cdot 0.02}{1.024} = 9.77\,\mu S\), i.e. \(\frac{1}{9.77\,\mu S} = 102\,k\Omega\),
    which is the \(Z_{out} \approx 100\,k\Omega\) above (the \(2\%\) difference is the
    \(1 + K_{CLM} V_{DS}\) factor dropped in \(g_{o} \approx K_{CLM} I_D\)). Sizing \(W_0\) with
    \(K_{CLM} = 0\) also leaves \(I_{M_0}\) high by the factor \(1.024\) at \(V_{ov_0} = 0.2V\);
    for exactly \(100\,\mu A\), divide \(W_0\) and \(W_1\) by \(1.024\).
  \item \textbf{Compute the maximum load that can be connected
    at the output:}
    \begin{align*}
      V_{OUT,\min} &= V_{ee} + V_{ov_1} = V_{ee} + V_{ov_0} = 0 + 0.2V = 0.2V \\
      R_{L, \max} &= \frac{V_{cc} - V_{OUT,\min}}{I_{M_1}} \\
                   &= \frac{5V - 0.2V}{500 \mu A} \\
                   &= 9.6 k\Omega
    \end{align*}
\end{enumerate}

\subsubsection{Code}
\begin{lstlisting}[language=Octave]
  clear; clc;
  pkg load symbolic;

  syms Vx Vout Iin Itest
  syms G3 G2 go0 gm0 go1 gm1

  % small signal kcl, equations (5) and (6)
  eq5 = Vout*(G3 + go1) + gm1*Vx == 0;   % KCL at Vout
  eq6 = Vx*(G2 + go0 + gm0) == 0;        % KCL at Vx

  % voltage transfer: Vx forced by a source, so (6) drops out
  sol = solve(eq5, Vout);
  printf('=== Vout/Vx ===\n'); pretty(simplify(sol/Vx))

  % current gain: Iin into Vx, output shorted (Vout = 0)
  vx = solve(Vx*(G2 + go0 + gm0) == Iin, Vx);
  Iout = gm1*vx + go1*0;                 % M1 drain current
  printf('=== Iout/Iin ===\n'); pretty(simplify(Iout/Iin))

  % input impedance: Itest into Vx
  vx = solve(Vx*(G2 + go0 + gm0) == Itest, Vx);
  printf('=== Zin ===\n'); pretty(simplify(vx/Itest))

  % output impedance: Itest into Vout, (6) keeps Vx = 0
  sol = solve([Vout*(G3 + go1) + gm1*Vx == Itest, eq6], [Vout, Vx]);
  printf('=== Zout ===\n'); pretty(simplify(sol.Vout/Itest))
\end{lstlisting}

\subsection{Wilson Mirror}

\subsubsection{Overview}
Constructed from three transistors, it improves
upon the Basic Mirror by providing a higher output
impedance (negative feedback through \(M_0\) holds
\(V_X\), and so \(I_{M_1} = I_{OUT}\), constant), which
results in better current source/sink performance.

\subsubsection{Schematic}
\begin{circuitfig}
  % Paths, nodes and wires:
  \node[shape=rectangle, draw={rgb,255:red,255;green,0;blue,0}, line width=1pt, dash pattern={on 4pt off 2pt on 1pt off 2pt on 1pt off 2pt}, minimum width=1.465cm, minimum height=2.195cm] at (6.25, 7.135){};
  \draw (3.23, 8.04) to[american voltage source, l_={$V_{cc}$}] (1.98, 8.04);
  \node[nmos, arrowmos, xscale=-1] (N1) at (3.5, 2.75){} node[anchor=east] at (N1.text){$M_0$};
  \node[nmos, arrowmos] (N2) at (6, 2.75){} node[anchor=west] at (N2.text){$M_1$};
  \draw (3.5, 7.54) to[european resistor, l={$G_3$}] (3.5, 6.04);
  \draw (6, 5.79) |- (6, 5.25);
  \draw (4.75, 2.75) |- (6, 3.71);
  \draw (5.02, 2.75) -- (4.48, 2.75);
  \draw (3.5, 1.98) -| (3.5, 1.75) -| (6, 1.98);
  \draw (4, 1) to[american voltage source, l={$V_{ee}$}] (2.75, 1);
  \draw (4, 1) -| (4.75, 1.75);
  \node[ground, rotate=-90] at (1.98, 8.04){};
  \node[ground, rotate=-90] at (2.75, 1){};
  \draw[dash pattern={on 0.4pt off 0.8pt}, -stealth] (6, 5.79) -- (6.5, 5.79);
  \node[shape=rectangle, minimum width=0.965cm, minimum height=0.465cm] (N3) at (7, 5.75){} node[anchor=center] at (N3.text){$V_{out}$};
  \draw (6, 7.5) to[european resistor, l={\textcolor{rgb,255:red,255;green,0;blue,0}{$G_4$}}] (6, 6);
  \draw (6, 5.79) -- (6, 6.02);
  \draw (3.5, 7.54) |- (3.23, 8.04);
  \draw (3.5, 8.04) -| (6, 7.5);
  \draw (10.25, 10.75) -- (10.25, -0);
  \node[shape=rectangle, minimum width=4.59cm, minimum height=1.09cm] at (5.187, 9.813){} node[anchor=north, align=center, text width=4.202cm, inner sep=6pt] at (5.187, 10.375){Large Signal Schematic\\(Sink)};
  \node[shape=rectangle, minimum width=3.09cm, minimum height=2.465cm] at (8.687, 7){} node[anchor=north west, align=left, text width=2.702cm, inner sep=6pt] at (7.125, 8.25){\textcolor{rgb,255:red,255;green,0;blue,0}{\footnotesize Exists Only For Analysis. In  Implementation take output directly from V\_\{out\}}};
  \node[shape=rectangle, minimum width=0.965cm, minimum height=0.465cm] (N4) at (7, 3.71){} node[anchor=center] at (N4.text){$V_{x}$};
  \draw[dash pattern={on 0.4pt off 0.8pt}, -stealth] (6, 3.71) |- (6.5, 3.71);
  \draw (3.5, 4.5) -| (3.5, 6.04);
  \draw (17.625, 2.562) to[american current source, l_={$g_{m_1}V_{gs}$}] (17.625, 1.062);
  \draw (19.125, 2.562) to[european resistor, l_={$g_{o_1}$}] (19.125, 1.062);
  \draw (17.625, 1.062) -- (19.125, 1.062);
  \draw (19.125, 2.562) -- (17.625, 2.562);
  \draw (17.625, 1.062) -- (16.375, 1.062);
  \draw (16.375, 2.562) -- (15.625, 2.562);
  \node[shape=rectangle, minimum width=0.465cm, minimum height=0.465cm] at (16.375, 2.312){} node[anchor=center, align=center, text width=0.077cm, inner sep=6pt] at (16.375, 2.312){\footnotesize G};
  \node[shape=rectangle, minimum width=0.465cm, minimum height=0.465cm] at (16.375, 1.312){} node[anchor=center, align=center, text width=0.077cm, inner sep=6pt] at (16.375, 1.312){\footnotesize S};
  \node[shape=rectangle, minimum width=0.465cm, minimum height=0.465cm] at (19.625, 2.312){} node[anchor=center, align=center, text width=0.077cm, inner sep=6pt] at (19.625, 2.312){\footnotesize D};
  \draw (15.625, 2.562) -- (14.875, 2.562);
  \draw (13.375, 2.562) to[american current source, l={$g_{m_0}V_{gs}$}] (13.375, 1.062);
  \draw (11.875, 2.562) to[european resistor, l={$g_{o_0}$}] (11.875, 1.062);
  \draw (13.375, 2.563) -- (11.875, 2.563);
  \draw (13.375, 1.063) -- (11.875, 1.063);
  \draw (14.875, 1.062) -- (13.375, 1.062);
  \node[shape=rectangle, minimum width=0.465cm, minimum height=0.465cm] at (14.875, 2.312){} node[anchor=center, align=center, text width=0.077cm, inner sep=6pt] at (14.875, 2.312){\footnotesize G};
  \node[shape=rectangle, minimum width=0.465cm, minimum height=0.465cm] at (14.875, 1.312){} node[anchor=center, align=center, text width=0.077cm, inner sep=6pt] at (14.875, 1.312){\footnotesize S};
  \node[shape=rectangle, minimum width=0.465cm, minimum height=0.465cm] at (11.375, 2.312){} node[anchor=center, align=center, text width=0.077cm, inner sep=6pt] at (11.375, 2.312){\footnotesize D};
  \draw (11.875, 7.563) to[european resistor, l={$G_3$}] (11.875, 6.063);
  \draw (11.875, 6.063) -| (11.875, 5.562);
  \draw (11.875, 7.563) -| (11.875, 8.062) -- (12.625, 8.062) -| (12.625, 7.812);
  \draw (19.125, 7.563) to[european resistor, l={\textcolor{rgb,255:red,255;green,0;blue,0}{$G_4$}}] (19.125, 6.063);
  \draw (19.125, 6.063) -| (19.125, 5.562);
  \draw (19.125, 7.563) -| (19.125, 8.062) |- (18.375, 8.062) -| (18.375, 7.812);
  \node[ground] at (14.875, 1.062){};
  \node[ground] at (16.375, 1.062){};
  \node[ground] at (12.625, 7.812){};
  \node[ground] at (18.375, 7.812){};
  \node[shape=rectangle, minimum width=4.215cm, minimum height=1.09cm] at (15.5, 9.813){} node[anchor=north, align=center, text width=3.827cm, inner sep=6pt] at (15.5, 10.375){Small Signal Schematic\\(Sink)};
  \draw (6, 3.52) -| (6, 3.71);
  \node[nmos, arrowmos] (N5) at (6, 4.48){} node[anchor=west] at (N5.text){$M_2$};
  \draw (5.02, 4.48) -| (3.5, 3.52);
  \node[shape=rectangle, minimum width=0.965cm, minimum height=0.465cm] (N6) at (2.5, 4.5){} node[anchor=center] at (N6.text){$V_{y}$};
  \draw[dash pattern={on 0.4pt off 0.8pt}, -stealth] (3.5, 4.48) -- (3, 4.48);
  \draw (17.625, 5.562) to[american current source, l_={$g_{m_2}V_{gs}$}] (17.625, 4.062);
  \draw (19.125, 5.562) to[european resistor, l_={$g_{o_2}$}] (19.125, 4.062);
  \draw (17.625, 4.062) -- (19.125, 4.062);
  \draw (19.125, 5.562) -- (17.625, 5.562);
  \draw (17.625, 4.062) -- (16.375, 4.062);
  \draw (16.375, 5.562) -- (15.625, 5.562);
  \node[shape=rectangle, minimum width=0.465cm, minimum height=0.465cm] at (16.375, 5.312){} node[anchor=center, align=center, text width=0.077cm, inner sep=6pt] at (16.375, 5.312){\footnotesize G};
  \node[shape=rectangle, minimum width=0.465cm, minimum height=0.465cm] at (16.375, 4.312){} node[anchor=center, align=center, text width=0.077cm, inner sep=6pt] at (16.375, 4.312){\footnotesize S};
  \draw (19.125, 4.062) -- (19.125, 2.562);
  \draw (15.625, 2.562) -| (15.625, 3.313) -- (19.125, 3.313);
  \node[circ] at (3.5, 8.04){};
  \node[circ] at (3.5, 4.48){};
  \node[circ] at (4.75, 2.75){};
  \node[circ] at (6, 3.71){};
  \node[circ] at (4.75, 1.75){};
  \node[circ] at (15.625, 2.562){};
  \node[circ] at (19.125, 3.313){};
  \draw (11.875, 2.562) |- (15.625, 5.562);
  \node[shape=rectangle, minimum width=0.965cm, minimum height=0.465cm] (N7) at (10.875, 5.582){} node[anchor=center] at (N7.text){$V_{y}$};
  \draw[dash pattern={on 0.4pt off 0.8pt}, -stealth] (11.875, 5.562) |- (11.375, 5.562);
  \node[circ] at (11.875, 5.562){};
  \draw[dash pattern={on 0.4pt off 0.8pt}, -stealth] (19.125, 5.562) -- (19.625, 5.562);
  \node[shape=rectangle, minimum width=0.965cm, minimum height=0.465cm] (N8) at (20.125, 5.522){} node[anchor=center] at (N8.text){$V_{out}$};
  \node[shape=rectangle, minimum width=0.965cm, minimum height=0.465cm] (N9) at (14.625, 3.333){} node[anchor=center] at (N9.text){$I_G$};
  \node[shape=rectangle, minimum width=0.965cm, minimum height=0.465cm] (N10) at (20.125, 3.313){} node[anchor=center] at (N10.text){$V_{x}$};
  \draw[dash pattern={on 0.4pt off 0.8pt}, -stealth] (19.125, 3.313) |- (19.625, 3.313);
  \node[shape=rectangle, minimum width=0.965cm, minimum height=0.465cm] (N11) at (3.75, 3.73){} node[anchor=center] at (N11.text){$I_G$};
  \node[shape=rectangle, draw={rgb,255:red,255;green,0;blue,0}, line width=1pt, dash pattern={on 4pt off 2pt on 1pt off 2pt on 1pt off 2pt}, minimum width=1.59cm, minimum height=2.527cm] at (19.312, 7.094){};
  \draw[dash pattern={on 0.4pt off 0.8pt}, -stealth] (4.75, 3.71) -- (4.25, 3.71);
  \draw[dash pattern={on 0.4pt off 0.8pt}, -stealth] (15.625, 3.313) -- (15.125, 3.313);
\end{circuitfig}

\subsubsection{Transfer Function}
Here's a table of the voltages at each terminal of the transistors:

\begin{tabularx}{\linewidth}{|X|X|X|X|X|X|} \hline
  Transistor & Gate (G) & Source (S) & Drain (D) & $V_{GS}$ & $V_{DS}$ \\ \hline
  $M_0$ & $V_{X}$ & $V_{ee}$ & $V_{Y}$ & $V_{X} - V_{ee}$ & $V_{Y} - V_{ee}$ \\ \hline
  $M_1$ & $V_{X}$ & $V_{ee}$ & $V_{X}$ & $V_{X} - V_{ee}$ & $V_{X} - V_{ee}$ \\ \hline
  $M_2$ & $V_{Y}$ & $V_{X}$ & $V_{OUT}$ & $V_{Y} - V_{X}$ & $V_{OUT} - V_{X}$ \\ \hline
\end{tabularx}

\begin{enumerate}
  \item \textbf{Large Signal, Time Domain, DC Input Analysis (Q5)}:
    Applying KCL at $V_{OUT}$:
    \begin{align*}
      (V_{OUT} - V_{cc})G_4 + I_{M_2} &= 0 \hspace{1cm} \text{[KCL at } V_{OUT}\text{]} \\
      V_{OUT} &= V_{cc} - \frac{I_{M_2}}{G_4} \tag{1} \\
      (V_{Y} - V_{cc})G_3 + I_{M_0} + I_{G_2} &= 0 \hspace{1cm} \text{[KCL at } V_{Y}\text{, } I_{G_2} \text{ into the gate of } M_2\text{]} \\
      (V_{Y} - V_{cc})G_3 + I_{M_0} &= 0 \hspace{1cm} \text{[} I_{G_2} \approx 0 \text{ by MOS physics]} \\
      V_{Y} &= V_{cc} - \frac{I_{M_0}}{G_3} \tag{2} \\
      I_{M_1} + I_{G} - I_{M_2} &= 0 \hspace{1cm} \text{[KCL at } V_{X}\text{]} \\
      I_{M_1} + I_{G} &= I_{M_2} \tag{3}
    \end{align*}
    Using transistor equations for $I_{M_0}$, $I_{M_1}$ and $I_{M_2}$
    and applying the voltages:
    \begin{align*}
      I_{M_0} &= k_{n_0}(V_{GS_0} - V_{th})^2 \cdot (1 + K_{CLM} V_{DS_0}) \\
              &= k_{n_0}((V_{X} - V_{ee}) - V_{th})^2 \cdot (1 + K_{CLM} (V_{Y} - V_{ee})) \\
              &= k_{n_0}\left(V_{X} - V_{ee} - V_{th}\right)^2 \cdot \left(1 + K_{CLM} \left(V_{cc} - \frac{I_{M_0}}{G_3} - V_{ee}\right)\right) \\
              &= k_{n_0}\left(V_{X} - V_{ee} - V_{th}\right)^2 \quad \text{[assuming } K_{CLM} = 0 \text{]} \tag{4} \\ \\
      I_{M_1} &= k_{n_1}(V_{GS_1} - V_{th})^2 \cdot (1 + K_{CLM} V_{DS_1}) \\
              &= k_{n_1}((V_{X} - V_{ee}) - V_{th})^2 \cdot (1 + K_{CLM} (V_{X} - V_{ee})) \\
              &= k_{n_1}\left(V_{X} - V_{ee} - V_{th}\right)^2 \cdot \left(1 + K_{CLM} (V_{X} - V_{ee})\right) \\
              &= k_{n_1}\left(V_{X} - V_{ee} - V_{th}\right)^2 \quad \text{[assuming } K_{CLM} = 0 \text{]} \tag{5} \\ \\
      I_{M_2} &= k_{n_2}(V_{GS_2} - V_{th})^2 \cdot (1 + K_{CLM} V_{DS_2}) \\
              &= k_{n_2}((V_{Y} - V_{X}) - V_{th})^2 \cdot (1 + K_{CLM} (V_{OUT} - V_{X})) \\
              &= k_{n_2}\left(V_{Y} - V_{X} - V_{th}\right)^2 \cdot \left(1 + K_{CLM} \left(V_{cc} - \frac{I_{M_2}}{G_4} - V_{X}\right)\right) \\
              &= k_{n_2}\left(V_{Y} - V_{X} - V_{th}\right)^2 \quad \text{[assuming } K_{CLM} = 0 \text{]} \tag{6}
    \end{align*}
    The input current is the current through \(G_3\) and the output
    current is the drain current of \(M_2\):
    \begin{align*}
      I_{IN} &= (V_{cc} - V_Y)G_3 = I_{M_0} \quad \text{[from (2)]} \\
      I_{OUT} &= I_{M_2} = I_{M_1} + I_{G} \approx I_{M_1} \quad \text{[from (3), } I_{G} \approx 0 \text{ by MOS physics]}
    \end{align*}
    \(M_0\) and \(M_1\) share \(V_{GS} = V_X - V_{ee}\), so dividing the
    second line of (5) by the second line of (4) cancels the square term:
    \begin{align*}
      \frac{I_{M_1}}{I_{M_0}} &= \frac{k_{n_1}(V_{X} - V_{ee} - V_{th})^2 (1 + K_{CLM} (V_{X} - V_{ee}))}{k_{n_0}(V_{X} - V_{ee} - V_{th})^2 (1 + K_{CLM} (V_{Y} - V_{ee}))} \\
                              &= \frac{k_{n_1}}{k_{n_0}} \cdot \frac{1 + K_{CLM} (V_{X} - V_{ee})}{1 + K_{CLM} (V_{Y} - V_{ee})} \\
                              &= \frac{k_{n_1}}{k_{n_0}} \quad \text{[assuming } K_{CLM} = 0 \text{]}
    \end{align*}
    With \(K_{CLM} \neq 0\) the factor stays, since
    \(V_Y - V_X = V_{GS_2} \neq 0\): \(V_{DS_0} \neq V_{DS_1}\). This is the
    systematic error removed by the Modified Wilson Mirror. Thus, with
    \(k_n = \frac{K_P}{2}\frac{W}{L}\),
    \begin{align*}
      \frac{I_{OUT}}{I_{IN}} \approx \frac{I_{M_2}}{I_{M_0}} \approx \frac{I_{M_1}}{I_{M_0}} \approx \frac{k_{n_1}}{k_{n_0}} = \frac{\frac{K_P}{2} \frac{W_1}{L_1}}{\frac{K_P}{2} \frac{W_0}{L_0}} = \frac{\frac{W_1}{L_1}}{\frac{W_0}{L_0}} = \frac{W_1 L_0}{L_1 W_0}
    \end{align*}
  \item \textbf{Small Signal, Frequency Domain, AC Input Analysis (Q4)}:\\
    Applying KCL at $V_{out}$:
    \begin{align*}
      V_{out}G_4 + (V_{out} - V_{x})g_{o_2} + g_{m_2}(V_{y} - V_{x}) &= 0 \\
      V_{out}(G_4 + g_{o_2}) - V_{x}(g_{o_2} + g_{m_2}) + g_{m_2}V_{y} &= 0 \tag{7}
    \end{align*}
    Applying KCL at $V_{x}$:
    \begin{align*}
      (V_{x} - V_{out})g_{o_2} + V_{x}g_{o_1} + g_{m_1}V_{x} - g_{m_2}(V_{y} - V_{x}) + I_{G} &= 0 \\
      (V_{x} - V_{out})g_{o_2} + V_{x}g_{o_1} + g_{m_1}V_{x} - g_{m_2}(V_{y} - V_{x}) &= 0 \quad \text{[as } I_{G} \approx 0 \text{ by MOS physics]} \\
      V_{x}g_{o_2} - V_{out}g_{o_2} + V_{x}g_{o_1} + g_{m_1}V_{x} - g_{m_2}V_{y} + g_{m_2}V_{x} &= 0 \\
      V_{x}(g_{o_2} + g_{o_1} + g_{m_1} + g_{m_2}) - V_{out}g_{o_2} - g_{m_2}V_{y} &= 0 \tag{8}
    \end{align*}
    Applying KCL at $V_{y}$:
    \begin{align*}
      V_{y}G_3 + V_{y}g_{o_0} + g_{m_0}V_{x} &= 0 \\
      V_{y}(G_3 + g_{o_0}) + g_{m_0}V_{x} &= 0 \tag{9}
    \end{align*}
    Adding voltage source $V_s$ at $V_y$ that supplies
    a current $I_s$ into $V_y$, we get:
    \begin{align*}
      V_{y}(G_3 + g_{o_0}) + g_{m_0}V_{x} - I_{s} &= 0 \tag{10}
    \end{align*}
    Fourth equation arises from Voltage source:
    \begin{align*}
      V_{y} &= V_{s} \tag{11}
    \end{align*}
    Solving equations (7), (8), (10) and (11),
    \begin{align*}
      V_{out}(G_4 + g_{o_2}) - V_{x}(g_{o_2} + g_{m_2}) + g_{m_2}V_{y} &= 0 \\
      V_{out}(G_4 + g_{o_2}) - V_{x}(g_{o_2} + g_{m_2}) + g_{m_2}V_{s} &= 0 \quad \text{[from (11)]} \tag{12} \\
      \\
      V_{x}(g_{o_2} + g_{o_1} + g_{m_1} + g_{m_2}) - V_{out}g_{o_2} - g_{m_2}V_{y} &= 0 \\
      V_{x}(g_{o_2} + g_{o_1} + g_{m_1} + g_{m_2}) - V_{out}g_{o_2} - g_{m_2}V_{s} &= 0 \quad \text{[from (11)]} \\
      V_{x}(g_{o_2} + g_{o_1} + g_{m_1} + g_{m_2}) &= V_{out}g_{o_2} + g_{m_2}V_{s} \\
      V_{x} &= \frac{V_{out}g_{o_2} + g_{m_2}V_{s}}{g_{o_2} + g_{o_1} + g_{m_1} + g_{m_2}} \tag{13}
    \end{align*}
    Putting (13) in (12),
    \begin{align*}
      V_{out}(G_4 + g_{o_2}) - \left(\frac{V_{out}g_{o_2} + g_{m_2}V_{s}}{g_{o_2} + g_{o_1} + g_{m_1} + g_{m_2}}\right)(g_{o_2} + g_{m_2}) + g_{m_2}V_{s} &= 0 \\
      V_{out}(G_4 + g_{o_2}) - \frac{V_{out}g_{o_2}(g_{o_2} + g_{m_2})}{g_{o_2} + g_{o_1} + g_{m_1} + g_{m_2}} - \frac{g_{m_2}V_{s}(g_{o_2} + g_{m_2})}{g_{o_2} + g_{o_1} + g_{m_1} + g_{m_2}} + g_{m_2}V_{s} &= 0 \\
      V_{out}\left((G_4 + g_{o_2}) - \frac{g_{o_2}(g_{o_2} + g_{m_2})}{g_{o_2} + g_{o_1} + g_{m_1} + g_{m_2}}\right) + V_{s}\left(g_{m_2} - \frac{g_{m_2}(g_{o_2} + g_{m_2})}{g_{o_2} + g_{o_1} + g_{m_1} + g_{m_2}}\right) &= 0 \\
      V_{out}\left(\frac{(G_4 + g_{o_2})(g_{o_2} + g_{o_1} + g_{m_1} + g_{m_2}) - g_{o_2}(g_{o_2} + g_{m_2})}{g_{o_2} + g_{o_1} + g_{m_1} + g_{m_2}}\right) \\ \qquad \, = -V_{s}\left(\frac{g_{m_2}(g_{o_2} + g_{o_1} + g_{m_1} + g_{m_2}) - g_{m_2}(g_{o_2} + g_{m_2})}{g_{o_2} + g_{o_1} + g_{m_1} + g_{m_2}}\right) \\
      V_{out}\left(\frac{G_4(g_{o_2} + g_{o_1} + g_{m_1} + g_{m_2}) + g_{o_2}(g_{o_2} + g_{o_1} + g_{m_1} + g_{m_2}) - g_{o_2}(g_{o_2} + g_{m_2})}{g_{o_2} + g_{o_1} + g_{m_1} + g_{m_2}}\right) \\ \qquad \, = -V_{s}\left(\frac{g_{m_2}(g_{o_2} + g_{o_1} + g_{m_1} + g_{m_2}) - g_{m_2}(g_{o_2} + g_{m_2})}{g_{o_2} + g_{o_1} + g_{m_1} + g_{m_2}}\right) \\
      V_{out}\left(\frac{G_4(g_{o_2} + g_{o_1} + g_{m_1} + g_{m_2}) + g_{o_2}(g_{o_1} + g_{m_1}) + g_{o_2}(g_{o_2} + g_{m_2}) - g_{o_2}(g_{o_2} + g_{m_2})}{g_{o_2} + g_{o_1} + g_{m_1} + g_{m_2}}\right) \\ \qquad \, = -V_{s}\left(\frac{g_{m_2}(g_{o_2} + g_{o_1} + g_{m_1} + g_{m_2}) - g_{m_2}(g_{o_2} + g_{m_2})}{g_{o_2} + g_{o_1} + g_{m_1} + g_{m_2}}\right) \\
      V_{out}\left(\frac{G_4(g_{o_2} + g_{o_1} + g_{m_1} + g_{m_2}) + g_{o_2}(g_{o_1} + g_{m_1})}{g_{o_2} + g_{o_1} + g_{m_1} + g_{m_2}}\right) \\ \qquad \, = -V_{s}\left(\frac{g_{m_2}(g_{o_2} + g_{o_1} + g_{m_1} + g_{m_2}) - g_{m_2}(g_{o_2} + g_{m_2})}{g_{o_2} + g_{o_1} + g_{m_1} + g_{m_2}}\right) \\
      V_{out}(G_4(g_{o_2} + g_{o_1} + g_{m_1} + g_{m_2}) + g_{o_2}(g_{o_1} + g_{m_1})) \\ \qquad \, = -V_{s}(g_{m_2}(g_{o_2} + g_{o_1} + g_{m_1} + g_{m_2}) - g_{m_2}(g_{o_2} + g_{m_2})) \\\\
      V_{out}(G_4(g_{o_2} + g_{o_1} + g_{m_1} + g_{m_2}) + g_{o_2}(g_{o_1} + g_{m_1})) \\ \qquad \, = -V_{s}(g_{m_2}g_{o_2} + g_{m_2}g_{o_1} + g_{m_2}g_{m_1} + g_{m_2}^2 - g_{m_2}g_{o_2} - g_{m_2}^2) \\\\
      V_{out}(G_4(g_{o_2} + g_{o_1} + g_{m_1} + g_{m_2}) + g_{o_2}(g_{o_1} + g_{m_1})) \\ \qquad \, = -V_{s}(g_{m_2}g_{o_1} + g_{m_2}g_{m_1}) \\\\
      V_{out}(G_4(g_{o_2} + g_{o_1} + g_{m_1} + g_{m_2}) + g_{o_2}(g_{o_1} + g_{m_1})) \\ \qquad \, = -V_{s}(g_{m_2}(g_{o_1} + g_{m_1})) \\\\
      \therefore \frac{V_{out}}{V_s} = \frac{-g_{m_2}(g_{o_1} + g_{m_1})}{G_4(g_{o_2} + g_{o_1} + g_{m_1} + g_{m_2}) + g_{o_2}(g_{o_1} + g_{m_1})} \tag{14}
    \end{align*}
\end{enumerate}

\subsubsection{Analysis}
\begin{enumerate}
  \item \textbf{Impedances, Small Signal, Frequency Domain, AC Input (Q4)}:
    \paragraph{At Input:} Applying a test current $I_{test}$
      at $V_{y}$, from equation (9):
      \begin{align*}
        V_{y}(G_3 + g_{o_0}) + g_{m_0}V_{x} &= I_{test}
      \end{align*}
      Substituting (14) into (13) (\(V_s = V_y\) by (11), and (14) holds for
      any \(V_y\) since it follows from (7) and (8) only):
      \begin{align*}
        V_{x} &= \frac{V_{out}g_{o_2} + g_{m_2}V_{y}}{g_{o_2} + g_{o_1} + g_{m_1} + g_{m_2}} \\
              &= \frac{\left(V_{y} \cdot \frac{-g_{m_2}(g_{o_1} + g_{m_1})}{G_4(g_{o_2} + g_{o_1} + g_{m_1} + g_{m_2}) + g_{o_2}(g_{o_1} + g_{m_1})}\right)g_{o_2} + g_{m_2}V_{y}}{g_{o_2} + g_{o_1} + g_{m_1} + g_{m_2}} \\
              &= \frac{V_{y}\frac{-g_{m_2}g_{o_2}(g_{o_1} + g_{m_1})}{G_4(g_{o_2} + g_{o_1} + g_{m_1} + g_{m_2}) + g_{o_2}(g_{o_1} + g_{m_1})} + g_{m_2}V_{y}}{g_{o_2} + g_{o_1} + g_{m_1} + g_{m_2}} \\
              &= V_{y} \cdot \frac{g_{m_2} - \frac{g_{m_2}g_{o_2}(g_{o_1} + g_{m_1})}{G_4(g_{o_2} + g_{o_1} + g_{m_1} + g_{m_2}) + g_{o_2}(g_{o_1} + g_{m_1})}}{g_{o_2} + g_{o_1} + g_{m_1} + g_{m_2}} \\
              &= V_{y} \cdot \frac{g_{m_2}G_4(g_{o_2} + g_{o_1} + g_{m_1} + g_{m_2}) + g_{m_2}g_{o_2}(g_{o_1} + g_{m_1}) - g_{m_2}g_{o_2}(g_{o_1} + g_{m_1})}{(g_{o_2} + g_{o_1} + g_{m_1} + g_{m_2})G_4(g_{o_2} + g_{o_1} + g_{m_1} + g_{m_2}) + g_{o_2}(g_{o_1} + g_{m_1})(g_{o_2} + g_{o_1} + g_{m_1} + g_{m_2})} \\
              &= V_{y} \cdot \frac{g_{m_2}G_4(g_{o_2} + g_{o_1} + g_{m_1} + g_{m_2})}{(g_{o_2} + g_{o_1} + g_{m_1} + g_{m_2})G_4(g_{o_2} + g_{o_1} + g_{m_1} + g_{m_2}) + g_{o_2}(g_{o_1} + g_{m_1})(g_{o_2} + g_{o_1} + g_{m_1} + g_{m_2})} \\
              &= V_{y} \cdot \frac{g_{m_2}G_4}{G_4(g_{o_2} + g_{o_1} + g_{m_1} + g_{m_2}) + g_{o_2}(g_{o_1} + g_{m_1})} \tag{15}
      \end{align*}
      Substituting (15) into (9):
      \begin{align*}
        V_{y}(G_3 + g_{o_0}) + g_{m_0}\left(V_{y} \cdot \frac{g_{m_2}G_4}{G_4(g_{o_2} + g_{o_1} + g_{m_1} + g_{m_2}) + g_{o_2}(g_{o_1} + g_{m_1})}\right) &= I_{test} \\
        V_{y}\left((G_3 + g_{o_0}) + \frac{g_{m_0}g_{m_2}G_4}{G_4(g_{o_2} + g_{o_1} + g_{m_1} + g_{m_2}) + g_{o_2}(g_{o_1} + g_{m_1})}\right) &= I_{test} \\
        V_{y}\left(\frac{(G_3 + g_{o_0})(G_4(g_{o_2} + g_{o_1} + g_{m_1} + g_{m_2}) + g_{o_2}(g_{o_1} + g_{m_1})) + g_{m_0}g_{m_2}G_4}{G_4(g_{o_2} + g_{o_1} + g_{m_1} + g_{m_2}) + g_{o_2}(g_{o_1} + g_{m_1})}\right) &= I_{test} \\
        \therefore Z_{in} = \frac{V_{y}}{I_{test}} = \frac{G_4(g_{o_2} + g_{o_1} + g_{m_1} + g_{m_2}) + g_{o_2}(g_{o_1} + g_{m_1})}{(G_3 + g_{o_0})(G_4(g_{o_2} + g_{o_1} + g_{m_1} + g_{m_2}) + g_{o_2}(g_{o_1} + g_{m_1})) + g_{m_0}g_{m_2}G_4}
      \end{align*}
    \paragraph{At Output:} Applying a test current $I_{test}$
      at $V_{out}$ (the input current source is opened, so (9) holds
      unchanged). Let \(S = g_{o_1} + g_{m_1} + g_{o_2} + g_{m_2}\) and
      \(H = G_3 + g_{o_0}\). From (7), (8) and (9):
      \begin{align*}
        V_{out}(G_4 + g_{o_2}) - V_{x}(g_{o_2} + g_{m_2}) + g_{m_2}V_{y} &= I_{test} \tag{16} \\
        V_{x} S - V_{out}g_{o_2} - g_{m_2}V_{y} &= 0 \tag{17} \\
        V_{y} H + g_{m_0}V_{x} = 0 \quad \Rightarrow \quad V_{y} &= -\frac{g_{m_0}}{H} V_{x} \tag{18}
      \end{align*}
      Substituting (18) into (17):
      \begin{align*}
        V_{x} S + \frac{g_{m_0} g_{m_2}}{H} V_{x} &= V_{out}g_{o_2} \\
        V_{x} \left( \frac{S H + g_{m_0} g_{m_2}}{H} \right) &= V_{out}g_{o_2} \\
        V_{x} &= \frac{g_{o_2} H}{S H + g_{m_0} g_{m_2}} V_{out} \tag{19}
      \end{align*}
      Substituting (18) into (16):
      \begin{align*}
        V_{out}(G_4 + g_{o_2}) - V_{x}(g_{o_2} + g_{m_2}) - \frac{g_{m_0} g_{m_2}}{H} V_{x} &= I_{test} \\
        V_{out}(G_4 + g_{o_2}) - V_{x} \left( \frac{(g_{o_2} + g_{m_2}) H + g_{m_0} g_{m_2}}{H} \right) &= I_{test} \tag{20}
      \end{align*}
      Substituting (19) into (20):
      \begin{align*}
        V_{out}(G_4 + g_{o_2}) - V_{out} \frac{g_{o_2} \left( (g_{o_2} + g_{m_2}) H + g_{m_0} g_{m_2} \right)}{S H + g_{m_0} g_{m_2}} &= I_{test} \\
        V_{out} \left( \frac{(G_4 + g_{o_2})(S H + g_{m_0} g_{m_2}) - g_{o_2}(g_{o_2} + g_{m_2}) H - g_{o_2} g_{m_0} g_{m_2}}{S H + g_{m_0} g_{m_2}} \right) &= I_{test} \\
        V_{out} \left( \frac{G_4 (S H + g_{m_0} g_{m_2}) + g_{o_2} S H - g_{o_2}(g_{o_2} + g_{m_2}) H}{S H + g_{m_0} g_{m_2}} \right) &= I_{test} \\
        V_{out} \left( \frac{G_4 (S H + g_{m_0} g_{m_2}) + g_{o_2} H (g_{o_1} + g_{m_1})}{S H + g_{m_0} g_{m_2}} \right) &= I_{test} \hspace{1cm} [S - g_{o_2} - g_{m_2} = g_{o_1} + g_{m_1}]
      \end{align*}
      \begin{align*}
        \therefore Z_{out} = \frac{V_{out}}{I_{test}} = \frac{S H + g_{m_0} g_{m_2}}{G_4 (S H + g_{m_0} g_{m_2}) + g_{o_2} H (g_{o_1} + g_{m_1})}
      \end{align*}
      which is the load \(G_4\) in parallel with the mirror's own output
      impedance. Dividing numerator and denominator by \(S H + g_{m_0} g_{m_2}\):
      \begin{align*}
        Z_{out} &= \frac{1}{G_4 + \frac{g_{o_2} H (g_{o_1} + g_{m_1})}{S H + g_{m_0} g_{m_2}}} = \frac{1}{G_4 + \frac{1}{Z_{out,mirror}}}
      \end{align*}
      so the mirror alone is (\(G_4 \to 0\)):
      \begin{align*}
        Z_{out,mirror} &= \frac{S H + g_{m_0} g_{m_2}}{g_{o_2} (g_{o_1} + g_{m_1}) H} = \frac{(G_3 + g_{o_0}) S + g_{m_0} g_{m_2}}{g_{o_2} (g_{o_1} + g_{m_1}) (G_3 + g_{o_0})}
      \end{align*}
      With an ideal input current source (\(G_3 \to 0\), \(H = g_{o_0}\)) and
      \(g_m \gg g_o\) (\(S \approx g_{m_1} + g_{m_2}\), \(g_{o_1} + g_{m_1} \approx g_{m_1}\)):
      \begin{align*}
        Z_{out,mirror} &= \frac{g_{o_0} S + g_{m_0} g_{m_2}}{g_{o_2} (g_{o_1} + g_{m_1}) g_{o_0}} \\
                       &\approx \frac{g_{o_0}(g_{m_1} + g_{m_2}) + g_{m_0} g_{m_2}}{g_{o_2} g_{m_1} g_{o_0}} \\
                       &\approx \frac{g_{m_0} g_{m_2}}{g_{o_0} g_{o_2} g_{m_1}} \quad \text{[} g_{o_0}(g_{m_1} + g_{m_2}) \ll g_{m_0} g_{m_2} \text{]}
      \end{align*}
      With \(g_m = \frac{2I}{V_{ov}}\), \(g_o = K_{CLM} I\), equal \(V_{ov}\) and
      \(I_{M_1} = I_{M_2} = I_{OUT}\), \(I_{M_0} = I_{IN}\):
      \begin{align*}
        Z_{out,mirror} &\approx \frac{\frac{2I_{IN}}{V_{ov}} \cdot \frac{2I_{OUT}}{V_{ov}}}{K_{CLM} I_{IN} \cdot K_{CLM} I_{OUT} \cdot \frac{2I_{OUT}}{V_{ov}}}
                        = \frac{\frac{4 I_{IN} I_{OUT}}{V_{ov}^2}}{\frac{2 K_{CLM}^2 I_{IN} I_{OUT}^2}{V_{ov}}}
                        = \frac{4 I_{IN} I_{OUT} V_{ov}}{2 K_{CLM}^2 I_{IN} I_{OUT}^2 V_{ov}^2}
                        = \frac{2}{K_{CLM}^2 I_{OUT} V_{ov}}
      \end{align*}
  \item \textbf{Output Load (Max), Large Signal, Time Domain, DC Input (Q5)}:
    \(M_2\) saturated needs \(V_{DS_2} \geq V_{ov_2} = V_{GS_2} - V_{th}\):
    \begin{align*}
      V_{OUT} - V_X &\geq V_{Y} - V_{X} - V_{th} \\
      V_{OUT} &\geq V_Y - V_{th} \\
      V_{OUT,\min} &= V_Y - V_{th} \\
      R_{L,\max} &= \frac{V_{cc} - V_{OUT,\min}}{I_{OUT}}
    \end{align*}
\end{enumerate}

\subsubsection{Design}
Large Signal, Time Domain, DC Input Analysis (Q5) can be used to set the bias
currents and voltages in this circuit.
\begin{enumerate}
  \item Choose appropriate transistor dimensions
    ($W$ and $L$) for transistors $M_0$ and $M_1$
    to achieve the desired current gain.
    \begin{align*}
      \frac{I_{OUT}}{I_{IN}} &= \frac{W_1 L_0}{L_1 W_0}
    \end{align*}
  \item Choose a suitable overdrive voltage value
    ($V_{ov}$) for setting the bias voltages at
    nodes $V_{X}$ and $V_{Y}$.
    \begin{align*}
      V_{ov_0} &= V_{GS_0} - V_{th} \\
      V_{ov_0} &= (V_{X} - V_{ee}) - V_{th} \\
      V_{X} &= V_{ov_0} + V_{th} + V_{ee} \\
      \\
      V_{ov_2} &= V_{GS_2} - V_{th} \\
      V_{ov_2} &= (V_{Y} - V_{X}) - V_{th} \\
      V_{Y} &= V_{ov_2} + V_{th} + V_{X}
    \end{align*}
  \item Using the chosen $V_{ov}$, compute $W_0$
    to set the required bias current $I_{IN}$.
    \begin{align*}
      I_{IN} = I_{M_0} &= \frac{K_P}{2} \frac{W_0}{L_0} V_{ov_0}^2 \\
      \therefore W_0 &= \frac{2 I_{IN} L_0}{K_P V_{ov_0}^2}
    \end{align*}
    $W_1$ can then be computed using the current
    gain relation from step (a). \(M_2\) carries \(I_{M_2} = I_{OUT}\)
    at the chosen \(V_{ov_2}\):
    \begin{align*}
      I_{OUT} = I_{M_2} &= \frac{K_P}{2} \frac{W_2}{L_2} V_{ov_2}^2 \\
      \therefore W_2 &= \frac{2 I_{OUT} L_2}{K_P V_{ov_2}^2}
    \end{align*}
  \item Compute the required value of $G_3$ (or
    $R_3$) to ensure that the voltage at node
    $V_{Y}$ is as calculated in step (b).
    \begin{align*}
      (V_{Y} - V_{cc})G_3 + I_{M_0} &= 0 \\
      (V_{cc} - V_{Y})G_3 &= I_{M_0} \\
      \therefore G_3 &= \frac{I_{IN}}{V_{cc} - V_{Y}} \quad \text{[}I_{M_0} = I_{IN} \text{]} \\
      R_3 &= \frac{1}{G_3} = \frac{V_{cc} - V_{Y}}{I_{IN}}
    \end{align*}
\end{enumerate}

\subsubsection{Question}
Design a Wilson Current Mirror (Sink) that meets the
following requirements. The mirror must provide a current
gain at the output of 5. The input current must be
$I_{IN} = 100 \mu A$. The mirror must operate with
$V_{cc} = 5V$ and $V_{ee} = 0V$. Assume the following
transistor parameters: $V_{th} = 1V$, $K_{CLM} = 0.02 V^{-1}$,
$K_P = 200 \mu A/V^2$ and
$L = 1 \mu m$ for all transistors. Find the output impedance
and the maximum possible load that can be connected at the
output, with the load returned to $V_{cc}$.
\paragraph{Solution:}
\begin{enumerate}
  \item \textbf{Choose appropriate transistor dimensions:}
    \begin{align*}
      \frac{I_{OUT}}{I_{IN}} &= \frac{W_1 L_0}{L_1 W_0} = 5 \\
      \text{Assuming } L_0 &= L_1 = 1 \mu m \\
      \frac{W_1 \cdot 1 \mu m}{1 \mu m \cdot W_0} &= 5 \\
      \therefore W_1 &= 5 W_0 \tag{i}
    \end{align*}
  \item \textbf{Choose a suitable overdrive voltage and determine
    the resulting bias voltages:} Assuming $V_{ov_0} = V_{ov_2} = 0.2V$,
    \begin{align*}
      V_{X} &= V_{ov_0} + V_{th} + V_{ee} = 0.2V + 1V + 0V = 1.2V \\
      V_{Y} &= V_{ov_2} + V_{th} + V_{X} = 0.2V + 1V + 1.2V = 2.4V
    \end{align*}
  \item \textbf{Compute $W_0$ to set the required bias current:}
    For $I_{IN} = 100 \mu A$
    \begin{align*}
      W_0 &= \frac{2 I_{IN} L_0}{K_P V_{ov_0}^2} \\
          &= \frac{2 \cdot 100 \mu A \cdot 1 \mu m}{200 \mu A/V^2 \cdot (0.2V)^2} \\
          &= 25 \mu m \\
      W_1 &= 5 W_0 = 5 \cdot 25 \mu m = 125 \mu m \quad \text{[from (i)]}
    \end{align*}
    \(M_2\) carries \(I_{OUT} = 5 \cdot 100 \mu A = 500 \mu A\) at \(V_{ov_2} = 0.2V\):
    \begin{align*}
      W_2 &= \frac{2 I_{OUT} L_2}{K_P V_{ov_2}^2} = \frac{2 \cdot 500 \mu A \cdot 1 \mu m}{200 \mu A/V^2 \cdot (0.2V)^2} = 125 \mu m
    \end{align*}
  \item \textbf{Compute the required value of $G_3$ (or $R_3$):}
    \begin{align*}
      R_3 &= \frac{1}{G_3} = \frac{V_{cc} - V_{Y}}{I_{IN}} \\
          &= \frac{5V - 2.4V}{100 \mu A} = \frac{2.6V}{100 \mu A} = 26k \Omega
    \end{align*}
  \item \textbf{Compute the output impedance:} \(M_2\) is sized like \(M_1\)
    (\(V_{ov_2} = 0.2V\) at \(500\,\mu A\)). With \(g_m = \frac{2 I_D}{V_{ov}}\),
    \(g_o = K_{CLM} I_D\):
    \begin{align*}
      g_{m_0} &= \frac{2 \cdot 100 \mu A}{0.2V} = 1 mS \hspace{1cm}
      g_{m_1} = g_{m_2} = \frac{2 \cdot 500 \mu A}{0.2V} = 5 mS \\
      g_{o_0} &= 0.02 \cdot 100 \mu A = 2 \mu S \hspace{1cm}
      g_{o_1} = g_{o_2} = 0.02 \cdot 500 \mu A = 10 \mu S \\
      G_3 &= \frac{1}{26 k\Omega} = 38.46 \mu S \hspace{1cm}
      S = g_{o_1} + g_{m_1} + g_{o_2} + g_{m_2} = 10 \mu S + 5 mS + 10 \mu S + 5 mS = 10.02 mS
    \end{align*}
    Ideal input current source (\(G_3 \to 0\)):
    \begin{align*}
      Z_{out,mirror} &\approx \frac{2}{K_{CLM}^2 I_{OUT} V_{ov}} = \frac{2}{0.02^2 \cdot 500 \mu A \cdot 0.2V} = \frac{2}{4 \times 10^{-8} \, A} = 50 M\Omega
    \end{align*}
    Input current set by \(R_3\) (\(H = G_3 + g_{o_0} = 40.46\,\mu S\)):
    \begin{align*}
      Z_{out,mirror} &= \frac{H S + g_{m_0} g_{m_2}}{g_{o_2} (g_{o_1} + g_{m_1}) H} = \frac{(40.46 \mu)(10.02 m) + (1 m)(5 m)}{(10 \mu)(5.01 m)(40.46 \mu)} \\
                     &= \frac{0.405 \times 10^{-6} + 5 \times 10^{-6}}{2.027 \times 10^{-12}} = \frac{5.405 \times 10^{-6}}{2.027 \times 10^{-12}} \approx 2.67 M\Omega
    \end{align*}
    \(R_3\) shunts the feedback node \(V_Y\), which is why the Wilson mirror
    is fed from a high impedance current source in practice.
  \item \textbf{Compute the maximum load that can be connected
    at the output:}
    \begin{align*}
      V_{OUT,\min} &= V_Y - V_{th} = 2.4V - 1V = 1.4V \\
      R_{L, \max} &= \frac{V_{cc} - V_{OUT,\min}}{I_{M_2}} = \frac{5V - 1.4V}{500 \mu A} = 7.2 k\Omega
    \end{align*}
\end{enumerate}

\subsubsection{Code}
\begin{lstlisting}[language=Octave]
  clear; clc;
  pkg load symbolic;

  syms Vy Vx Vout Itest
  syms G3 G4 go0 gm0 go1 gm1 go2 gm2

  % small signal kcl, equations (7), (8) and (9)
  eq7 = Vout*(G4 + go2) - Vx*(go2 + gm2) + gm2*Vy == 0;          % KCL at Vout
  eq8 = Vx*(go2 + go1 + gm1 + gm2) - Vout*go2 - gm2*Vy == 0;     % KCL at Vx
  eq9 = Vy*(G3 + go0) + gm0*Vx == 0;                             % KCL at Vy

  % voltage transfer (14): Vy forced by a source, so (9) drops out
  sol = solve([eq7, eq8], [Vout, Vx]);
  printf('=== Vout/Vy ===\n'); pretty(factor(simplify(sol.Vout/Vy)))

  % input impedance: Itest into Vy
  sol = solve([eq7, eq8, Vy*(G3 + go0) + gm0*Vx == Itest], [Vout, Vx, Vy]);
  printf('=== Zin ===\n'); pretty(factor(simplify(sol.Vy/Itest)))

  % output impedance (16): Itest into Vout, load G4 kept
  eq16 = Vout*(G4 + go2) - Vx*(go2 + gm2) + gm2*Vy == Itest;
  sol = solve([eq16, eq8, eq9], [Vout, Vx, Vy]);
  zout = factor(simplify(sol.Vout/Itest));
  printf('=== Zout ===\n'); pretty(zout)
  printf('=== Zout, mirror alone (G4 = 0) ===\n'); pretty(factor(subs(zout, G4, 0)))
\end{lstlisting}

\subsection{Modified Wilson Mirror}

\subsubsection{Overview}
Constructed from four transistors, It improves
upon the Wilson Current Mirror by adding the diode
connected \(M_3\) so that \(V_{ds_0} = V_{ds_1}\)
(\(V_y = V_x\)). Both mirroring transistors then see the
same channel length modulation and the systematic
current error of the Wilson mirror disappears.

\subsubsection{Schematic}
\begin{circuitfig}
  % Paths, nodes and wires:
  \node[shape=rectangle, draw={rgb,255:red,255;green,0;blue,0}, line width=1pt, dash pattern={on 4pt off 2pt on 1pt off 2pt on 1pt off 2pt}, minimum width=1.715cm, minimum height=2.195cm] at (6.125, 7.135){};
  \draw (3.23, 8.04) to[american voltage source, l_={$V_{cc}$}] (1.98, 8.04);
  \node[nmos, arrowmos, xscale=-1] (N1) at (3.5, 2.75){} node[anchor=east] at (N1.text){$M_0$};
  \node[nmos, arrowmos] (N2) at (6, 2.75){} node[anchor=west] at (N2.text){$M_1$};
  \draw (3.5, 7.54) to[european resistor, l={$G_4$}] (3.5, 6.04);
  \draw (6, 5.79) |- (6, 5.25);
  \draw (4.75, 2.75) |- (6, 3.71);
  \draw (5.02, 2.75) -- (4.48, 2.75);
  \draw (3.5, 1.98) -| (3.5, 1.75) -| (6, 1.98);
  \draw (4, 1) to[american voltage source, l={$V_{ee}$}] (2.75, 1);
  \draw (4, 1) -| (4.75, 1.75);
  \node[ground, rotate=-90] at (1.98, 8.04){};
  \node[ground, rotate=-90] at (2.75, 1){};
  \draw[dash pattern={on 0.4pt off 0.8pt}, -stealth] (6, 5.79) -- (6.5, 5.79);
  \node[shape=rectangle, minimum width=0.965cm, minimum height=0.465cm] (N3) at (7, 5.75){} node[anchor=center] at (N3.text){$V_{out}$};
  \draw (6, 7.5) to[european resistor, l={\textcolor{rgb,255:red,255;green,0;blue,0}{$G_5$}}] (6, 6);
  \draw (6, 5.79) -- (6, 6.02);
  \draw (3.5, 7.54) |- (3.23, 8.04);
  \draw (3.5, 8.04) -| (6, 7.5);
  \draw (10.25, 10.75) -- (10.25, -0);
  \node[shape=rectangle, minimum width=5.59cm, minimum height=1.09cm] at (5.438, 9.937){} node[anchor=north, align=center, text width=5.202cm, inner sep=6pt] at (5.438, 10.5){Large Signal Schematic\\(Sink)};
  \node[shape=rectangle, minimum width=3.09cm, minimum height=2.465cm] at (8.687, 7){} node[anchor=north west, align=left, text width=2.702cm, inner sep=6pt] at (7.125, 8.25){\textcolor{rgb,255:red,255;green,0;blue,0}{\footnotesize Exists Only For Analysis. In  Implementation take output directly from V\_\{out\}}};
  \node[shape=rectangle, minimum width=0.965cm, minimum height=0.465cm] (N4) at (7, 3.71){} node[anchor=center] at (N4.text){$V_{x}$};
  \draw[dash pattern={on 0.4pt off 0.8pt}, -stealth] (6, 3.71) |- (6.5, 3.71);
  \node[shape=rectangle, minimum width=0.965cm, minimum height=0.465cm] (N5) at (4.75, 4.25){} node[anchor=center] at (N5.text){$I_{G_1}$};
  \draw (17.625, 2.75) to[american current source, l_={$g_{m_1}V_{gs}$}] (17.625, 1.25);
  \draw (19.125, 2.75) to[european resistor, l_={$g_{o_1}$}] (19.125, 1.25);
  \draw (17.625, 1.25) -- (19.125, 1.25);
  \draw (19.125, 2.75) -- (17.625, 2.75);
  \draw (17.625, 1.25) -- (16.375, 1.25);
  \draw (16.375, 2.75) -- (15.625, 2.75);
  \node[shape=rectangle, minimum width=0.465cm, minimum height=0.465cm] at (16.375, 2.5){} node[anchor=center, align=center, text width=0.077cm, inner sep=6pt] at (16.375, 2.5){\footnotesize G};
  \node[shape=rectangle, minimum width=0.465cm, minimum height=0.465cm] at (16.375, 1.5){} node[anchor=center, align=center, text width=0.077cm, inner sep=6pt] at (16.375, 1.5){\footnotesize S};
  \node[shape=rectangle, minimum width=0.465cm, minimum height=0.465cm] at (19.625, 2.5){} node[anchor=center, align=center, text width=0.077cm, inner sep=6pt] at (19.625, 2.5){\footnotesize D};
  \draw (15.625, 2.75) -- (14.875, 2.75);
  \draw (13.375, 2.75) to[american current source, l={$g_{m_0}V_{gs}$}] (13.375, 1.25);
  \draw (11.875, 2.75) to[european resistor, l={$g_{o_0}$}] (11.875, 1.25);
  \draw (13.375, 2.75) -- (11.875, 2.75);
  \draw (13.375, 1.25) -- (11.875, 1.25);
  \draw (11.875, 2.75) -| (11.875, 4.25);
  \draw (14.875, 1.25) -- (13.375, 1.25);
  \node[shape=rectangle, minimum width=0.465cm, minimum height=0.465cm] at (14.875, 2.5){} node[anchor=center, align=center, text width=0.077cm, inner sep=6pt] at (14.875, 2.5){\footnotesize G};
  \node[shape=rectangle, minimum width=0.465cm, minimum height=0.465cm] at (14.875, 1.5){} node[anchor=center, align=center, text width=0.077cm, inner sep=6pt] at (14.875, 1.5){\footnotesize S};
  \node[shape=rectangle, minimum width=0.465cm, minimum height=0.465cm] at (11.375, 2.5){} node[anchor=center, align=center, text width=0.077cm, inner sep=6pt] at (11.375, 2.5){\footnotesize D};
  \draw (11.875, 8.562) to[european resistor, l={$G_4$}] (11.875, 7.062);
  \draw (11.875, 7.062) -| (11.875, 6.562);
  \draw (11.875, 8.562) -| (11.875, 9.062) -- (12.625, 9.062) -| (12.625, 8.812);
  \draw (19.125, 7.75) to[european resistor, l={\textcolor{rgb,255:red,255;green,0;blue,0}{$G_5$}}] (19.125, 6.25);
  \draw (19.125, 6.25) -| (19.125, 5.75);
  \draw (19.125, 7.75) -| (19.125, 8.25) |- (18.375, 8.25) -| (18.375, 8);
  \node[shape=rectangle, minimum width=0.715cm, minimum height=0.465cm] (N6) at (20, 3.5){} node[anchor=center] at (N6.text){$V_{x}$};
  \node[shape=rectangle, minimum width=0.965cm, minimum height=0.465cm] (N7) at (20.125, 5.75){} node[anchor=center] at (N7.text){$V_{out}$};
  \node[ground] at (14.875, 1.25){};
  \node[ground] at (16.375, 1.25){};
  \node[ground] at (12.625, 8.812){};
  \node[ground] at (18.375, 8){};
  \node[shape=rectangle, minimum width=0.965cm, minimum height=0.465cm] (N8) at (14.625, 3.5){} node[anchor=center] at (N8.text){$I_{G_1}$};
  \node[shape=rectangle, minimum width=4.84cm, minimum height=1.09cm] at (15.25, 9.937){} node[anchor=north, align=center, text width=4.452cm, inner sep=6pt] at (15.25, 10.5){Small Signal Schematic\\(Sink)};
  \draw (6, 3.52) -| (6, 3.71);
  \node[nmos, arrowmos] (N9) at (6, 4.48){} node[anchor=west] at (N9.text){$M_2$};
  \node[shape=rectangle, minimum width=0.965cm, minimum height=0.465cm] (N10) at (2.5, 3.71){} node[anchor=center] at (N10.text){$V_{y}$};
  \draw[dash pattern={on 0.4pt off 0.8pt}, -stealth] (3.5, 3.71) |- (3, 3.71);
  \draw (17.625, 5.75) to[american current source, l_={$g_{m_2}V_{gs}$}] (17.625, 4.25);
  \draw (19.125, 5.75) to[european resistor, l_={$g_{o_2}$}] (19.125, 4.25);
  \draw (17.625, 4.25) -- (19.125, 4.25);
  \draw (19.125, 5.75) -- (17.625, 5.75);
  \draw (17.625, 4.25) -- (16.375, 4.25);
  \draw (16.375, 5.75) -- (15.625, 5.75);
  \node[shape=rectangle, minimum width=0.465cm, minimum height=0.465cm] at (16.375, 5.5){} node[anchor=center, align=center, text width=0.077cm, inner sep=6pt] at (16.375, 5.5){\footnotesize G};
  \node[shape=rectangle, minimum width=0.465cm, minimum height=0.465cm] at (16.375, 4.5){} node[anchor=center, align=center, text width=0.077cm, inner sep=6pt] at (16.375, 4.5){\footnotesize S};
  \draw (19.125, 4.25) -- (19.125, 2.75);
  \draw (15.625, 2.75) -| (15.625, 3.5) -- (19.125, 3.5);
  \draw[dash pattern={on 0.4pt off 0.8pt}, -stealth] (19.125, 3.5) -- (19.625, 3.5);
  \node[shape=rectangle, minimum width=0.715cm, minimum height=0.465cm] (N11) at (11, 3.563){} node[anchor=center] at (N11.text){$V_{y}$};
  \draw[dash pattern={on 0.4pt off 0.8pt}, -stealth] (11.875, 3.563) -- (11.375, 3.563);
  \node[nmos, arrowmos, xscale=-1] (N12) at (3.5, 4.48){} node[anchor=east] at (N12.text){$M_3$};
  \draw (3.5, 3.71) -| (3.5, 3.52);
  \draw (4.48, 4.48) -- (5.02, 4.48);
  \draw (4.75, 4.5) |- (3.5, 5.5);
  \draw (3.5, 5.25) -| (3.5, 5.5);
  \draw (3.5, 6.04) -| (3.5, 5.5);
  \node[shape=rectangle, minimum width=0.965cm, minimum height=0.465cm] (N13) at (2.5, 5.5){} node[anchor=center] at (N13.text){$V_{z}$};
  \draw[dash pattern={on 0.4pt off 0.8pt}, -stealth] (3.5, 5.5) -- (3, 5.5);
  \node[shape=rectangle, minimum width=0.965cm, minimum height=0.465cm] (N14) at (4.75, 6){} node[anchor=center] at (N14.text){$I_{G_2}$};
  \draw (15.625, 5.75) -- (14.875, 5.75);
  \draw (13.375, 5.75) to[american current source, l={$g_{m_3}V_{gs}$}] (13.375, 4.25);
  \draw (11.875, 5.75) to[european resistor, l={$g_{o_3}$}] (11.875, 4.25);
  \draw (13.375, 5.75) -- (11.875, 5.75);
  \draw (13.375, 4.25) -- (11.875, 4.25);
  \draw (14.875, 4.25) -- (13.375, 4.25);
  \node[shape=rectangle, minimum width=0.465cm, minimum height=0.465cm] at (14.875, 5.5){} node[anchor=center, align=center, text width=0.077cm, inner sep=6pt] at (14.875, 5.5){\footnotesize G};
  \node[shape=rectangle, minimum width=0.465cm, minimum height=0.465cm] at (14.875, 4.5){} node[anchor=center, align=center, text width=0.077cm, inner sep=6pt] at (14.875, 4.5){\footnotesize S};
  \node[shape=rectangle, minimum width=0.465cm, minimum height=0.465cm] at (11.375, 5.5){} node[anchor=center, align=center, text width=0.077cm, inner sep=6pt] at (11.375, 5.5){\footnotesize D};
  \draw (15.625, 5.75) |- (11.875, 6.562) |- (11.875, 5.75);
  \node[shape=rectangle, minimum width=0.715cm, minimum height=0.465cm] (N15) at (11, 6.562){} node[anchor=center] at (N15.text){$V_{z}$};
  \draw[dash pattern={on 0.4pt off 0.8pt}, -stealth] (11.875, 6.562) -- (11.375, 6.562);
  \node[shape=rectangle, minimum width=0.965cm, minimum height=0.465cm] (N16) at (16.625, 6.562){} node[anchor=center] at (N16.text){$I_{G_2}$};
  \draw[dash pattern={on 0.4pt off 0.8pt}, -stealth] (19.125, 5.75) -- (19.625, 5.75);
  \draw[dash pattern={on 0.4pt off 0.8pt}, -stealth] (4.75, 5.5) -| (4.75, 5.75);
  \draw[dash pattern={on 0.4pt off 0.8pt}, -stealth] (4.779, 3.721) -| (4.75, 4);
  \draw[dash pattern={on 0.4pt off 0.8pt}, -stealth] (15.625, 6.562) -- (16.125, 6.562);
  \draw[dash pattern={on 0.4pt off 0.8pt}, -stealth] (15.625, 3.5) -- (15.125, 3.5);
  \node[shape=rectangle, draw={rgb,255:red,255;green,0;blue,0}, line width=1pt, dash pattern={on 4pt off 2pt on 1pt off 2pt on 1pt off 2pt}, minimum width=1.215cm, minimum height=2.195cm] at (19.375, 7.385){};
\end{circuitfig}

\subsubsection{Transfer Function}
Here's a table of the voltages at each terminal of the
transistors:\\
\begin{tabularx}{\linewidth}{|X|X|X|X|X|X|} \hline
  Transistor & Gate (G) & Source (S) & Drain (D) & $V_{gs}$ & $V_{ds}$ \\ \hline
  $M_0$ & $V_{x}$ & $V_{ee}$ & $V_{y}$ & $V_{x} - V_{ee}$ & $V_{y} - V_{ee}$ \\ \hline
  $M_1$ & $V_{x}$ & $V_{ee}$ & $V_{x}$ & $V_{x} - V_{ee}$ & $V_{x} - V_{ee}$ \\ \hline
  $M_2$ & $V_{z}$ & $V_{x}$ & $V_{out}$ & $V_{z} - V_{x}$ & $V_{out} - V_{x}$ \\ \hline
  $M_3$ & $V_{z}$ & $V_{y}$ & $V_{z}$ & $V_{z} - V_{y}$ & $V_{z} - V_{y}$ \\ \hline
\end{tabularx}\\

\begin{enumerate}
  \item \textbf{Large Signal, Time Domain, DC Input Analysis (Q5)}:\\
    Applying KCL at $V_{out}$:
    \begin{align*}
      (V_{out} - V_{cc})G_5 + I_{M_2} &= 0 \\
      V_{out}G_5 - V_{cc}G_5 + I_{M_2} &= 0 \\
      V_{out} &= \frac{V_{cc}G_5 - I_{M_2}}{G_5} \tag{1}
    \end{align*}
    Applying KCL at $V_{x}$:
    \begin{align*}
      I_{M_1} + I_{G_1} - I_{M_2} &= 0 \\
      I_{M_1} + I_{G_1} &= I_{M_2} \\
      I_{M_1} &\approx I_{M_2} \quad \text{[assuming } I_{G_1} \approx 0 \text{ by MOS physics]} \tag{2}
    \end{align*}
    Applying KCL at $V_{y}$:
    \begin{align*}
      I_{M_0} - I_{M_3} &= 0 \\
      I_{M_0} &= I_{M_3} \tag{3}
    \end{align*}
    Applying KCL at $V_{z}$:
    \begin{align*}
      (V_{z} - V_{cc})G_4 + I_{M_3} + I_{G_2} &= 0 \\
      V_{z}G_4 - V_{cc}G_4 + (I_{M_3} + I_{G_2}) &= 0 \\
      V_{z}G_4 &= V_{cc}G_4 - (I_{M_3} + I_{G_2}) \\
      V_{z} &= V_{cc} - \frac{(I_{M_3} + I_{G_2})}{G_4} \tag{4}
    \end{align*}
    Using transistor equations for $I_{M_0}$, $I_{M_1}$, $I_{M_2}$ and $I_{M_3}$
    and applying the voltages:
    \begin{align*}
      I_{M_0} &= k_{n_0}(V_{gs_0} - V_{th})^2 \cdot (1 + K_{CLM} V_{ds_0}) \\
              &= k_{n_0}((V_{x} - V_{ee}) - V_{th})^2 \cdot (1 + K_{CLM} (V_{y} - V_{ee})) \\
              &= k_{n_0}(V_{x} - V_{ee} - V_{th})^2 \quad \text{[assuming } K_{CLM} = 0 \text{]} \tag{5} \\ \\
      I_{M_1} &= k_{n_1}(V_{gs_1} - V_{th})^2 \cdot (1 + K_{CLM} V_{ds_1}) \\
              &= k_{n_1}((V_{x} - V_{ee}) - V_{th})^2 \cdot (1 + K_{CLM} (V_{x} - V_{ee})) \\
              &= k_{n_1}(V_{x} - V_{ee} - V_{th})^2 \quad \text{[assuming } K_{CLM} = 0 \text{]} \tag{6} \\ \\
      I_{M_2} &= k_{n_2}(V_{gs_2} - V_{th})^2 \cdot (1 + K_{CLM} V_{ds_2}) \\
              &= k_{n_2}((V_{z} - V_{x}) - V_{th})^2 \cdot (1 + K_{CLM} (V_{out} - V_{x})) \\
              &= k_{n_2}\left(\left(V_{cc} - \frac{(I_{M_3} + I_{G_2})}{G_4}\right) - V_{x} - V_{th}\right)^2 \cdot \left(1 + K_{CLM}\left(\frac{V_{cc}G_5 - I_{M_2}}{G_5}\right) - V_{x}\right) \\
              &= k_{n_2}\left(\left(V_{cc} - \frac{(I_{M_3} + I_{G_2})}{G_4}\right) - V_{x} - V_{th}\right)^2 \quad \text{[assuming } K_{CLM} = 0 \text{]} \tag{7} \\ \\
      I_{M_3} &= k_{n_3}(V_{gs_3} - V_{th})^2 \cdot (1 + K_{CLM} V_{ds_3}) \\
              &= k_{n_3}((V_{z} - V_{y}) - V_{th})^2 \cdot (1 + K_{CLM} (V_{z} - V_{y})) \\
              &= k_{n_3}\left(\left(V_{cc} - \frac{(I_{M_3} + I_{G_2})}{G_4}\right) - V_{y} - V_{th}\right)^2 \cdot \left(1 + K_{CLM}\left(\left(V_{cc} - \frac{(I_{M_3} + I_{G_2})}{G_4}\right) - V_{y}\right)\right) \\
              &= k_{n_3}\left(\left(V_{cc} - \frac{(I_{M_3} + I_{G_2})}{G_4}\right) - V_{y} - V_{th}\right)^2 \quad \text{[assuming } K_{CLM} = 0 \text{]} \tag{8}
    \end{align*}
    It can be seen that,
    \begin{align*}
      I_{in} &= I_{M_3} = I_{M_0} \quad \text{[from (3)]} \\
      I_{out} &= I_{M_2} = I_{M_1} \quad \text{[from (2)]} \\
      I_{M_0} &= I_{M_1} \quad \text{[from (5) and (6) and assuming } k_{n_0} = k_{n_1} \text{]} \\
      \therefore  I_{in} &= I_{out}
    \end{align*}
    Thus,
    \begin{align*}
      \frac{I_{out}}{I_{in}} \approx \frac{I_{M_2}}{I_{M_3}} \approx \frac{I_{M_1}}{I_{M_0}} \approx \frac{k_{n_1}}{k_{n_0}} = \frac{\frac{K_P W_1}{L_1}}{\frac{K_P W_0}{L_0}} = \frac{\frac{W_1}{L_1}}{\frac{W_0}{L_0}} = \frac{W_1 L_0}{L_1 W_0} = \frac{W_2 L_3}{L_2 W_3}
    \end{align*}
  \item \textbf{Small Signal, Frequency Domain, AC Input Analysis (Q4)}:\\
    Applying KCL at $V_{out}$:
    \begin{align*}
      V_{out}G_5 + (V_{out} - V_{x})g_{o_2} + g_{m_2}(V_{z} - V_{x}) &= 0 \\
      V_{out}G_5 + V_{out}g_{o_2} - V_{x}g_{o_2} + g_{m_2}V_{z} - g_{m_2}V_{x} &= 0 \\
      V_{out}(G_5 + g_{o_2}) - V_{x}(g_{o_2} + g_{m_2}) + g_{m_2}V_{z} &= 0 \tag{9}
    \end{align*}
    Applying KCL at $V_{x}$:
    \begin{align*}
      (V_{x} - V_{out})g_{o_2} + V_{x}g_{o_1} + g_{m_1}V_{x} - g_{m_2}(V_{z} - V_{x}) &= 0 \\
      V_{x}g_{o_2} - V_{out}g_{o_2} + V_{x}g_{o_1} + g_{m_1}V_{x} - g_{m_2}V_{z} + g_{m_2}V_{x} &= 0 \\
      V_{x}(g_{o_2} + g_{o_1} + g_{m_1} + g_{m_2}) - V_{out}g_{o_2} - g_{m_2}V_{z} &= 0 \tag{10}
    \end{align*}
    Applying KCL at $V_{y}$:
    \begin{align*}
      (V_{y} - V_{z})g_{o_3} - g_{m_3}(V_{z} - V_{y}) + V_{y}g_{o_0} + g_{m_0}V_{x} &= 0 \\
      V_{y}g_{o_3} - V_{z}g_{o_3} - g_{m_3}V_{z} + g_{m_3}V_{y} + V_{y}g_{o_0} + g_{m_0}V_{x} &= 0 \\
      V_{y}(g_{o_3} + g_{m_3} + g_{o_0}) - V_{z}(g_{o_3} + g_{m_3}) + g_{m_0}V_{x} &= 0 \tag{11}
    \end{align*}
    Applying KCL at $V_{z}$:
    \begin{align*}
      V_{z}G_4 + (V_{z} - V_{y})g_{o_3} + g_{m_3}(V_{z} - V_{y}) &= 0 \\
      V_{z}G_4 + V_{z}g_{o_3} - V_{y}g_{o_3} + g_{m_3}V_{z} - g_{m_3}V_{y} &= 0 \\
      V_{z}(G_4 + g_{o_3} + g_{m_3}) - V_{y}(g_{o_3} + g_{m_3}) &= 0 \tag{12}
    \end{align*}
    Adding voltage source $V_s$ at $V_z$ that supplies
    a current $I_s$ into $V_z$, we get:
    \begin{align*}
      V_{z}(G_4 + g_{o_3} + g_{m_3}) - V_{y}(g_{o_3} + g_{m_3}) - I_s &= 0 \tag{13}
    \end{align*}
    Fifth equation arises from Voltage source:
    \begin{align*}
      V_{z} &= V_{s} \tag{14}
    \end{align*}
    Equations (9), (10), (11), (13) and (14) can now be
    solved analytically, but its cumbersome, so I'll give the
    solution directly from Octave:
    \begin{align*}
      \frac{V_{out}}{V_{s}} = \frac{-g_{m_2}(g_{m_1} + g_{o_1})}{G_5(g_{m_1} + g_{m_2} + g_{o_1} + g_{o_2}) + g_{o_2}(g_{m_1} + g_{o_1})} \tag{15}
    \end{align*}
\end{enumerate}

\subsubsection{Analysis}
\begin{enumerate}
  \item \textbf{Impedances}: Let \(S = g_{o_1} + g_{m_1} + g_{o_2} + g_{m_2}\),
    \(P = g_{m_3} + g_{o_3}\) and \(M = G_4(P + g_{o_0}) + g_{o_0} P\).
    From (12), \(V_{y} = \frac{G_4 + P}{P} V_{z}\) whenever no test source
    sits at \(V_z\). Substituting into (11):
    \begin{align*}
      V_{z}\frac{(G_4 + P)(P + g_{o_0})}{P} - V_{z} P + g_{m_0} V_{x} &= 0 \\
      V_{z}\frac{(G_4 + P)(P + g_{o_0}) - P^2}{P} &= -g_{m_0} V_{x} \\
      V_{z}\frac{G_4 (P + g_{o_0}) + g_{o_0} P}{P} &= -g_{m_0} V_{x} \\
      V_{z} &= -\frac{g_{m_0} P}{M} V_{x} \tag{16}
    \end{align*}
    \paragraph{At Output:} Apply \(I_{test}\) at \(V_{out}\) (\(G_5\) removed),
      so (9) becomes
      \begin{align*}
        V_{out} g_{o_2} - V_{x}(g_{o_2} + g_{m_2}) + g_{m_2}V_{z} &= I_{test} \tag{17}
      \end{align*}
      Adding (17) and (10):
      \begin{align*}
        V_{x} S - V_{x}(g_{o_2} + g_{m_2}) &= I_{test} \\
        V_{x} (g_{o_1} + g_{m_1}) &= I_{test} \tag{18}
      \end{align*}
      Substituting (16) into (10):
      \begin{align*}
        V_{x} S + \frac{g_{m_0} g_{m_2} P}{M} V_{x} &= V_{out} g_{o_2} \\
        V_{x} \frac{S M + g_{m_0} g_{m_2} P}{M} &= V_{out} g_{o_2} \tag{19}
      \end{align*}
      Substituting (18) into (19):
      \begin{align*}
        \frac{I_{test}}{g_{o_1} + g_{m_1}} \cdot \frac{S M + g_{m_0} g_{m_2} P}{M} &= V_{out} g_{o_2} \\
        \therefore Z_{out} = \frac{V_{out}}{I_{test}} &= \frac{S M + g_{m_0} g_{m_2} P}{g_{o_2} (g_{o_1} + g_{m_1}) M}
      \end{align*}
      With an ideal input current source (\(G_4 \to 0\), \(M = g_{o_0} P\))
      and \(g_m \gg g_o\):
      \begin{align*}
        Z_{out} &= \frac{S g_{o_0} P + g_{m_0} g_{m_2} P}{g_{o_2}(g_{o_1} + g_{m_1}) g_{o_0} P}
                \approx \frac{g_{m_0} g_{m_2}}{g_{o_0} g_{o_2} g_{m_1}} = \frac{2}{K_{CLM}^2 I_{out} V_{ov}}
      \end{align*}
      With a load \(G_5\), the total is \(\frac{1}{G_5 + 1/Z_{out}}\).
    \paragraph{At Input:} Apply \(I_{test}\) at \(V_{z}\) with the output at AC
      ground (\(V_{out} = 0\)). (10) gives \(V_{x} = \frac{g_{m_2}}{S} V_{z}\);
      (11) gives
      \begin{align*}
        V_{y}(P + g_{o_0}) &= V_{z} P - g_{m_0} \frac{g_{m_2}}{S} V_{z} \\
        V_{y} &= \frac{P S - g_{m_0} g_{m_2}}{S (P + g_{o_0})} V_{z}
      \end{align*}
      and (12) with the test current:
      \begin{align*}
        V_{z}(G_4 + P) - V_{y} P &= I_{test} \\
        V_{z}\left( (G_4 + P) - \frac{P(P S - g_{m_0} g_{m_2})}{S (P + g_{o_0})} \right) &= I_{test} \\
        V_{z}\left( \frac{S\left[(G_4 + P)(P + g_{o_0}) - P^2\right] + g_{m_0} g_{m_2} P}{S (P + g_{o_0})} \right) &= I_{test} \\
        \therefore Z_{in} = \frac{V_{z}}{I_{test}} &= \frac{S (P + g_{o_0})}{S M + g_{m_0} g_{m_2} P}
      \end{align*}
      With \(G_4 \to 0\) and \(g_m \gg g_o\):
      \(Z_{in} \approx \frac{S P}{g_{m_0} g_{m_2} P} = \frac{g_{m_1} + g_{m_2}}{g_{m_0} g_{m_2}}\),
      which is \(\frac{2}{g_{m_0}}\) for \(g_{m_1} = g_{m_2}\): a low input impedance,
      as a current input should have.
  \item \textbf{Output Load (Max)}: \(M_2\) saturated needs
    \(V_{out} \geq V_z - V_{th}\):
    \begin{align*}
      R_{L,\max} &= \frac{V_{cc} - (V_z - V_{th})}{I_{out}}
    \end{align*}
\end{enumerate}

\subsubsection{Design}
Large Signal, Time Domain, DC Input Analysis (Q5) can be used to set the bias
currents and voltages in this circuit.
\begin{enumerate}
  \item Choose appropriate transistor dimensions
    ($W$ and $L$) for transistors $M_0$ and $M_1$
    to achieve the desired current gain.
    \begin{align*}
      \frac{I_{out}}{I_{in}} &= \frac{W_1 L_0}{L_1 W_0} = \frac{W_2 L_3}{L_2 W_3} \\
    \end{align*}
  \item Choose a suitable overdrive voltage value
    ($V_{ov}$) for computing the bias voltages at
    nodes $V_{x}$, $V_{y}$, and $V_{z}$.
    \begin{align*}
      V_{ov_0} &= V_{gs_0} - V_{th} \\
      V_{ov_0} &= V_{x} - V_{ee} - V_{th} \\
      V_{x} &= V_{ov_0} + V_{ee} + V_{th} \\
      \\
      V_{ov_2} &= V_{gs_2} - V_{th} \\
      V_{ov_2} &= V_{z} - V_{x} - V_{th} \\
      V_{z} &= V_{ov_2} + V_{x} + V_{th} \\
      \\
      V_{ov_3} &= V_{gs_3} - V_{th} \\
      V_{ov_3} &= V_{z} - V_{y} - V_{th} \\
      V_{y} &= V_{z} - V_{ov_3} - V_{th}
    \end{align*}
  \item Using the chosen $V_{ov}$, compute $W_0$
    to set the required bias current $I_{in}$.
    \begin{align*}
      I_{in} &\approx I_{M_0} = \frac{K_P}{2} \frac{W_0}{L_0} V_{ov_0}^2 \\
             &\approx I_{M_3} = \frac{K_P}{2} \frac{W_3}{L_3} V_{ov_3}^2 \\
      \therefore W_0 &= \frac{2 I_{in} L_0}{K_P V_{ov_0}^2} \\
                 W_3 &= \frac{2 I_{in} L_3}{K_P V_{ov_3}^2}
    \end{align*}
    $W_1$ can then be computed using the current
    gain relation from step (a). Similarly, $W_2$
    can be computed using the current gain relation
    from step (a).
  \item Compute the required value of $G_4$ (or
    $R_4$) to ensure that the voltage at node
    $V_{z}$ is as calculated in step (b).
    \begin{align*}
      V_{z} &= V_{cc} - \frac{(I_{M_3} + I_{G_2})}{G_4} \\
      V_{z} &\approx V_{cc} - \frac{I_{M_3}}{G_4} \quad \text{[assuming } I_{G_2} \approx 0 \text{ by MOS physics]} \\
      V_{z} - V_{cc} &\approx - \frac{I_{M_3}}{G_4} \\
      \frac{(I_{M_3})}{(V_{cc} - V_{z})} &\approx G_4 \\
      \frac{(I_{in})}{(V_{cc} - V_{z})} &\approx G_4 \quad \text{[} I_{in} = I_{M_3} = I_{M_0} \text{]} \\
      \therefore R_4 &\approx \frac{V_{cc} - V_{z}}{I_{in}}
    \end{align*}
\end{enumerate}

\subsubsection{Question}
Design a Modified Wilson Current Mirror in
Sink configuration to achieve a current gain of
2, with an input bias current of $I_{in} =
10\mu A$. Assume the following parameters for
the NMOS transistors used:
\begin{align*}
  K_P = 100 \mu A/V^2 \spc V_{th} = 0.7 V \spc K_{CLM} = 0.02 V^{-1} \\
  V_{cc} = 5 V \spc V_{ee} = 0 V \spc L_0 = L_1 = L_2 = L_3 = 1 \mu m
\end{align*}
Find the output impedance and the maximum
possible load that can be connected at the output
assuming $V_{L,\max} = V_{cc}$
\paragraph{Solution:}
\begin{enumerate}
  \item \textbf{Choose appropriate transistor dimensions:}
    \begin{align*}
      \frac{I_{out}}{I_{in}} &= \frac{W_1 L_0}{L_1 W_0} \\
      2 &= \frac{W_1 \cdot 1 \mu m}{1 \mu m \cdot W_0} \\
      W_1 &= 2 \cdot W_0
    \end{align*}
  \item \textbf{Choose a suitable overdrive voltage and determine
    the resulting bias voltages:}
    \begin{align*}
      V_{ov_0} &= 0.2 V \\
      V_{x} &= V_{ov_0} + V_{ee} + V_{th} = 0.2 + 0 + 0.7 = 0.9 V \\
      \\
      V_{ov_2} &= 0.2 V \\
      V_{z} &= V_{ov_2} + V_{x} + V_{th} = 0.2 + 0.9 + 0.7 = 1.8 V \\
      \\
      V_{ov_3} &= 0.2 V \\
      V_{y} &= V_{z} - V_{ov_3} - V_{th} = 1.8 - 0.2 - 0.7 = 0.9 V
    \end{align*}
  \item \textbf{Compute $W_0$ to set the required bias current:} For $I_{in} = 10 \mu A$
    \begin{align*}
      W_0 &= \frac{2 I_{in} L_0}{K_P V_{ov_0}^2} \\
          &= \frac{2 \cdot 10 \mu A \cdot 1 \mu m}{100 \mu A \cdot (0.2)^2} \\
          &= 5 \mu m
    \end{align*}
    Using (i) and (ii):
    \begin{align*}
      W_1 &= 2 \cdot W_0 = 2 \cdot 5 \mu m = 10 \mu m
    \end{align*}
  \item \textbf{Compute the required value of $G_4$ (or $R_4$):}
    \begin{align*}
      R_4 &\approx \frac{V_{cc} - V_{z}}{I_{in}} \\
          &= \frac{5 V - 1.8 V}{10 \mu A} \\
          &= 320 k\Omega
    \end{align*}
  \item \textbf{Output Load Limit:}
    \begin{align*}
      R_{L, \max} &= \frac{V_{cc} - (V_z - V_{th})}{I_{M_2}} = \frac{5 V - (1.8 V - 0.7 V)}{20 \mu A} = 195 k\Omega
    \end{align*}
  \item \textbf{Impedances:} \(I_{M_0} = I_{M_3} = 10\,\mu A\),
    \(I_{M_1} = I_{M_2} = 20\,\mu A\), all \(V_{ov} = 0.2\,V\):
    \begin{align*}
      g_{m_0} &= g_{m_3} = 100 \mu S \hspace{1cm} g_{m_1} = g_{m_2} = 200 \mu S \\
      g_{o_0} &= g_{o_3} = 0.2 \mu S \hspace{1cm} g_{o_1} = g_{o_2} = 0.4 \mu S \hspace{1cm} G_4 = \frac{1}{320 k\Omega} = 3.125 \mu S \\
      S &= 400.8 \mu S \hspace{1cm} P = 100.2 \mu S \hspace{1cm} M = G_4 (P + g_{o_0}) + g_{o_0} P = 333.8 \times 10^{-12}
    \end{align*}
    \paragraph{At Input:}
    \begin{align*}
      Z_{in} &= \frac{S (P + g_{o_0})}{S M + g_{m_0} g_{m_2} P} = \frac{(400.8 \mu)(100.4 \mu)}{(400.8 \mu)(333.8 p) + (100 \mu)(200 \mu)(100.2 \mu)} \approx 18.8 k\Omega
    \end{align*}
    \paragraph{At Output:}
    \begin{align*}
      Z_{out} &= \frac{S M + g_{m_0} g_{m_2} P}{g_{o_2} (g_{o_1} + g_{m_1}) M} \approx 79.9 M\Omega \hspace{1cm} [\text{with } R_4] \\
      Z_{out} &\approx \frac{2}{K_{CLM}^2 I_{out} V_{ov}} = \frac{2}{0.02^2 \cdot 20 \mu A \cdot 0.2 V} = 1.25 G\Omega \hspace{1cm} [\text{ideal input source}]
    \end{align*}
\end{enumerate}

\subsubsection{Code}
\begin{lstlisting}[language=Octave]
  clear; clc;
  pkg load symbolic;

  syms Vx Vy Vz Vout Vs Is Itest
  syms G5 G4 go0 gm0 go1 gm1 go2 gm2 go3 gm3

  eq1 = Vout*(G5 + go2) - Vx*(go2 + gm2) + gm2*Vz == 0; % KCL at Vout
  eq2 = Vx*(go2 + go1 + gm1 + gm2) - Vout*go2 - Vz*gm2 == 0; % KCL at Vx
  eq3 = Vy*(go3 + gm3 + go0) - Vz*(go3 + gm3) + Vx*gm0 == 0; % KCL at Vy
  eq4 = Vz*(G4 + go3 + gm3) - Vy*(go3 + gm3) == 0; % KCL at Vz

  % voltage transfer function: Vs forced at Vz
  eq4_vt = Vz*(G4 + go3 + gm3) - Vy*(go3 + gm3) - Is == 0;
  eq5_vt = Vz == Vs;
  sol_vt = solve([eq1, eq2, eq3, eq4_vt, eq5_vt], [Vout, Vx, Vy, Vz, Vs]);
  volt_transfer = simplify(sol_vt.Vout/sol_vt.Vs);
  % volt_transfer = collect(volt_transfer, G5);
  % volt_transfer = subs(volt_transfer, 'G4', 0);
  % volt_transfer = subs(volt_transfer, 'G5', 0);
  printf('=== Voltage Transfer Results ===\n');
  pretty(volt_transfer)
  printf('\n');

  % current transfer function: Is injected into Vz
  sol_ct = solve([eq1, eq2, eq3, eq4_vt], [Vout, Vx, Vy, Vz]);
  Iout = sol_ct.Vout * G5;
  curr_transfer = simplify(Iout/Is);
  % curr_transfer = collect(curr_transfer, G4*G5);
  printf('=== Current Transfer Results ===\n');
  pretty(curr_transfer)
  printf('\n');

  % input impedance: Itest injected into Vz
  eq4_zin = Vz*(G4 + go3 + gm3) - Vy*(go3 + gm3) - Itest == 0;
  sol_zin = solve([eq1, eq2, eq3, eq4_zin], [Vout, Vx, Vy, Vz]);
  zin = simplify(sol_zin.Vz/Itest);
  % zin = collect(zin, G4*G5);
  % zin = subs(zin, 'G4', 0);
  % zin = subs(zin, 'G5', 0);
  printf('=== Input Impedance Results ===\n');
  pretty(zin);
  printf('\n');

  % output impedance: Itest injected into Vout
  eq1_zout = Vout*(G5 + go2) - Vx*(go2 + gm2) + gm2*Vz - Itest == 0;
  sol_zout = solve([eq1_zout, eq2, eq3, eq4], [Vout, Vx, Vy, Vz]);
  zout = simplify(sol_zout.Vout/Itest);
  % zout = collect(zout, G4*G5);
  % zout = subs(zout, 'G4', 0);
  % zout = subs(zout, 'G5', 0);
  printf('=== Output Impedance Results ===\n');
  pretty(zout)
  printf('\n');
\end{lstlisting}

\subsection{Cascode Mirror}

\subsubsection{Overview}
Constructed from four transistors, It improves
output resistance and accuracy by reducing the
effects of channel length modulation.

\subsubsection{Schematic}
\begin{circuitfig}
  % Paths, nodes and wires:
  \node[shape=rectangle, draw={rgb,255:red,255;green,0;blue,0}, line width=1pt, dash pattern={on 4pt off 2pt on 1pt off 2pt on 1pt off 2pt}, minimum width=1.715cm, minimum height=2.195cm] at (6.125, 7.135){};
  \draw (3.23, 8.04) to[american voltage source, l_={$V_{cc}$}] (1.98, 8.04);
  \node[nmos, arrowmos, xscale=-1] (N1) at (3.5, 2.75){} node[anchor=east] at (N1.text){$M_0$};
  \node[nmos, arrowmos] (N2) at (6, 2.75){} node[anchor=west] at (N2.text){$M_1$};
  \draw (3.5, 7.54) to[european resistor, l={$G_4$}] (3.5, 6.04);
  \draw (6, 5.79) |- (6, 5.25);
  \draw (5.02, 2.75) -- (4.48, 2.75);
  \draw (3.5, 1.98) -| (3.5, 1.75) -| (6, 1.98);
  \draw (4, 1) to[american voltage source, l={$V_{ee}$}] (2.75, 1);
  \draw (4, 1) -| (4.75, 1.75);
  \node[ground, rotate=-90] at (1.98, 8.04){};
  \node[ground, rotate=-90] at (2.75, 1){};
  \draw[dash pattern={on 0.4pt off 0.8pt}, -stealth] (6, 5.5) -- (6.5, 5.5);
  \node[shape=rectangle, minimum width=0.965cm, minimum height=0.465cm] (N3) at (7, 5.46){} node[anchor=center] at (N3.text){$V_{out}$};
  \draw (6, 7.5) to[european resistor, l={\textcolor{rgb,255:red,255;green,0;blue,0}{$G_5$}}] (6, 6);
  \draw (6, 5.79) -- (6, 6.02);
  \draw (3.5, 7.54) |- (3.23, 8.04);
  \draw (3.5, 8.04) -| (6, 7.5);
  \draw (10.25, 10.75) -- (10.25, -0);
  \node[shape=rectangle, minimum width=5.59cm, minimum height=1.09cm] at (5.438, 9.937){} node[anchor=north, align=center, text width=5.202cm, inner sep=6pt] at (5.438, 10.5){Large Signal Schematic\\(Sink)};
  \node[shape=rectangle, minimum width=3.09cm, minimum height=2.465cm] at (8.687, 7){} node[anchor=north west, align=left, text width=2.702cm, inner sep=6pt] at (7.125, 8.25){\textcolor{rgb,255:red,255;green,0;blue,0}{\footnotesize Exists Only For Analysis. In  Implementation take output directly from V\_\{out\}}};
  \node[shape=rectangle, minimum width=0.965cm, minimum height=0.465cm] (N4) at (7, 3.71){} node[anchor=center] at (N4.text){$V_{y}$};
  \draw[dash pattern={on 0.4pt off 0.8pt}, -stealth] (6, 3.71) |- (6.5, 3.71);
  \node[shape=rectangle, minimum width=0.965cm, minimum height=0.465cm] (N5) at (4.75, 4.25){} node[anchor=center] at (N5.text){$I_{G_1}$};
  \draw (17.625, 2.75) to[american current source, l_={$g_{m_1}V_{gs}$}] (17.625, 1.25);
  \draw (19.125, 2.75) to[european resistor, l_={$g_{o_1}$}] (19.125, 1.25);
  \draw (17.625, 1.25) -- (19.125, 1.25);
  \draw (19.125, 2.75) -- (17.625, 2.75);
  \draw (17.625, 1.25) -- (16.375, 1.25);
  \draw (16.375, 2.75) -- (15.625, 2.75);
  \node[shape=rectangle, minimum width=0.465cm, minimum height=0.465cm] at (16.375, 2.5){} node[anchor=center, align=center, text width=0.077cm, inner sep=6pt] at (16.375, 2.5){\footnotesize G};
  \node[shape=rectangle, minimum width=0.465cm, minimum height=0.465cm] at (16.375, 1.5){} node[anchor=center, align=center, text width=0.077cm, inner sep=6pt] at (16.375, 1.5){\footnotesize S};
  \node[shape=rectangle, minimum width=0.465cm, minimum height=0.465cm] at (19.625, 2.5){} node[anchor=center, align=center, text width=0.077cm, inner sep=6pt] at (19.625, 2.5){\footnotesize D};
  \draw (15.625, 2.75) -- (14.875, 2.75);
  \draw (13.375, 2.75) to[american current source, l={$g_{m_0}V_{gs}$}] (13.375, 1.25);
  \draw (11.875, 2.75) to[european resistor, l={$g_{o_0}$}] (11.875, 1.25);
  \draw (13.375, 2.75) -- (11.875, 2.75);
  \draw (13.375, 1.25) -- (11.875, 1.25);
  \draw (11.875, 2.75) -| (11.875, 4.25);
  \draw (14.875, 1.25) -- (13.375, 1.25);
  \node[shape=rectangle, minimum width=0.465cm, minimum height=0.465cm] at (14.875, 2.5){} node[anchor=center, align=center, text width=0.077cm, inner sep=6pt] at (14.875, 2.5){\footnotesize G};
  \node[shape=rectangle, minimum width=0.465cm, minimum height=0.465cm] at (14.875, 1.5){} node[anchor=center, align=center, text width=0.077cm, inner sep=6pt] at (14.875, 1.5){\footnotesize S};
  \node[shape=rectangle, minimum width=0.465cm, minimum height=0.465cm] at (11.375, 2.5){} node[anchor=center, align=center, text width=0.077cm, inner sep=6pt] at (11.375, 2.5){\footnotesize D};
  \draw (11.875, 8.562) to[european resistor, l={$G_4$}] (11.875, 7.062);
  \draw (11.875, 7.062) -| (11.875, 6.562);
  \draw (11.875, 8.562) -| (11.875, 9.062) -- (12.625, 9.062) -| (12.625, 8.812);
  \draw (19.125, 7.75) to[european resistor, l={\textcolor{rgb,255:red,255;green,0;blue,0}{$G_5$}}] (19.125, 6.25);
  \draw (19.125, 6.25) -| (19.125, 5.75);
  \draw (19.125, 7.75) -| (19.125, 8.25) |- (18.375, 8.25) -| (18.375, 8);
  \node[shape=rectangle, minimum width=0.715cm, minimum height=0.465cm] (N6) at (20, 3.5){} node[anchor=center] at (N6.text){$V_{y}$};
  \node[shape=rectangle, minimum width=0.965cm, minimum height=0.465cm] (N7) at (20.125, 5.75){} node[anchor=center] at (N7.text){$V_{out}$};
  \node[ground] at (14.875, 1.25){};
  \node[ground] at (16.375, 1.25){};
  \node[ground] at (12.625, 8.812){};
  \node[ground] at (18.375, 8){};
  \node[shape=rectangle, minimum width=0.965cm, minimum height=0.465cm] (N8) at (16.625, 3.563){} node[anchor=center] at (N8.text){$I_{G_1}$};
  \node[shape=rectangle, minimum width=4.84cm, minimum height=1.09cm] at (15.25, 9.937){} node[anchor=north, align=center, text width=4.452cm, inner sep=6pt] at (15.25, 10.5){Small Signal Schematic\\(Sink)};
  \draw (6, 3.52) -| (6, 3.71);
  \node[nmos, arrowmos] (N9) at (6, 4.48){} node[anchor=west] at (N9.text){$M_2$};
  \node[shape=rectangle, minimum width=0.965cm, minimum height=0.465cm] (N10) at (2.5, 3.71){} node[anchor=center] at (N10.text){$V_{x}$};
  \draw[dash pattern={on 0.4pt off 0.8pt}, -stealth] (3.5, 3.71) |- (3, 3.71);
  \draw (17.625, 5.75) to[american current source, l_={$g_{m_2}V_{gs}$}] (17.625, 4.25);
  \draw (19.125, 5.75) to[european resistor, l_={$g_{o_2}$}] (19.125, 4.25);
  \draw (17.625, 4.25) -- (19.125, 4.25);
  \draw (19.125, 5.75) -- (17.625, 5.75);
  \draw (17.625, 4.25) -- (16.375, 4.25);
  \draw (16.375, 5.75) -- (15.625, 5.75);
  \node[shape=rectangle, minimum width=0.465cm, minimum height=0.465cm] at (16.375, 5.5){} node[anchor=center, align=center, text width=0.077cm, inner sep=6pt] at (16.375, 5.5){\footnotesize G};
  \node[shape=rectangle, minimum width=0.465cm, minimum height=0.465cm] at (16.375, 4.5){} node[anchor=center, align=center, text width=0.077cm, inner sep=6pt] at (16.375, 4.5){\footnotesize S};
  \draw (19.125, 4.25) -- (19.125, 2.75);
  \draw[dash pattern={on 0.4pt off 0.8pt}, -stealth] (19.125, 3.5) -- (19.625, 3.5);
  \node[shape=rectangle, minimum width=0.715cm, minimum height=0.465cm] (N11) at (11, 3.563){} node[anchor=center] at (N11.text){$V_{x}$};
  \draw[dash pattern={on 0.4pt off 0.8pt}, -stealth] (11.875, 3.563) -- (11.375, 3.563);
  \node[nmos, arrowmos, xscale=-1] (N12) at (3.5, 4.48){} node[anchor=east] at (N12.text){$M_3$};
  \draw (3.5, 3.71) -| (3.5, 3.52);
  \draw (4.48, 4.48) -- (5.02, 4.48);
  \draw (4.75, 4.5) |- (3.5, 5.5);
  \draw (3.5, 5.25) -| (3.5, 5.5);
  \draw (3.5, 6.04) -| (3.5, 5.5);
  \node[shape=rectangle, minimum width=0.965cm, minimum height=0.465cm] (N13) at (2.5, 5.5){} node[anchor=center] at (N13.text){$V_{z}$};
  \draw[dash pattern={on 0.4pt off 0.8pt}, -stealth] (3.5, 5.5) |- (3, 5.5);
  \node[shape=rectangle, minimum width=0.965cm, minimum height=0.465cm] (N14) at (4.75, 6){} node[anchor=center] at (N14.text){$I_{G_2}$};
  \draw (15.625, 5.75) -- (14.875, 5.75);
  \draw (13.375, 5.75) to[american current source, l={$g_{m_3}V_{gs}$}] (13.375, 4.25);
  \draw (11.875, 5.75) to[european resistor, l={$g_{o_3}$}] (11.875, 4.25);
  \draw (13.375, 5.75) -- (11.875, 5.75);
  \draw (13.375, 4.25) -- (11.875, 4.25);
  \draw (14.875, 4.25) -- (13.375, 4.25);
  \node[shape=rectangle, minimum width=0.465cm, minimum height=0.465cm] at (14.875, 5.5){} node[anchor=center, align=center, text width=0.077cm, inner sep=6pt] at (14.875, 5.5){\footnotesize G};
  \node[shape=rectangle, minimum width=0.465cm, minimum height=0.465cm] at (14.875, 4.5){} node[anchor=center, align=center, text width=0.077cm, inner sep=6pt] at (14.875, 4.5){\footnotesize S};
  \node[shape=rectangle, minimum width=0.465cm, minimum height=0.465cm] at (11.375, 5.5){} node[anchor=center, align=center, text width=0.077cm, inner sep=6pt] at (11.375, 5.5){\footnotesize D};
  \draw (15.625, 5.75) |- (11.875, 6.562) |- (11.875, 5.75);
  \node[shape=rectangle, minimum width=0.715cm, minimum height=0.465cm] (N15) at (11, 6.562){} node[anchor=center] at (N15.text){$V_{z}$};
  \draw[dash pattern={on 0.4pt off 0.8pt}, -stealth] (11.875, 6.562) -- (11.375, 6.562);
  \node[shape=rectangle, minimum width=0.965cm, minimum height=0.465cm] (N16) at (16.625, 6.562){} node[anchor=center] at (N16.text){$I_{G_2}$};
  \draw[dash pattern={on 0.4pt off 0.8pt}, -stealth] (19.125, 5.75) -- (19.625, 5.75);
  \draw[dash pattern={on 0.4pt off 0.8pt}, -stealth] (4.75, 5.5) -- (4.75, 5.75);
  \draw[dash pattern={on 0.4pt off 0.8pt}, -stealth] (4.75, 3.71) -- (4.75, 4);
  \draw[dash pattern={on 0.4pt off 0.8pt}, -stealth] (15.625, 6.562) -- (16.125, 6.562);
  \draw[dash pattern={on 0.4pt off 0.8pt}, -stealth] (15.625, 3.563) -- (16.125, 3.563);
  \node[shape=rectangle, draw={rgb,255:red,255;green,0;blue,0}, line width=1pt, dash pattern={on 4pt off 2pt on 1pt off 2pt on 1pt off 2pt}, minimum width=1.215cm, minimum height=2.195cm] at (19.375, 7.385){};
  \draw (3.5, 3.71) -| (4.75, 2.75);
  \draw (15.625, 2.75) |- (11.875, 3.563);
\end{circuitfig}

\subsubsection{Transfer Function}
Here's a table of the voltages at each terminal of the transistors:

\begin{tabularx}{\linewidth}{|X|X|X|X|X|X|} \hline
  Transistor & Gate (G) & Source (S) & Drain (D) & $V_{GS}$ & $V_{DS}$ \\ \hline
  $M_0$ & $V_{X}$ & $V_{ee}$ & $V_{X}$ & $V_{X} - V_{ee}$ & $V_{X} - V_{ee}$ \\ \hline
  $M_1$ & $V_{X}$ & $V_{ee}$ & $V_{Y}$ & $V_{X} - V_{ee}$ & $V_{Y} - V_{ee}$ \\ \hline
  $M_2$ & $V_{Z}$ & $V_{Y}$ & $V_{OUT}$ & $V_{Z} - V_{Y}$ & $V_{OUT} - V_{Y}$ \\ \hline
  $M_3$ & $V_{Z}$ & $V_{X}$ & $V_{Z}$ & $V_{Z} - V_{X}$ & $V_{Z} - V_{X}$ \\ \hline
\end{tabularx}

\begin{enumerate}
  \item \textbf{Large Signal, Time Domain, DC Input Analysis (Q5)}:\\
    Applying KCL at $V_{OUT}$:
    \begin{align*}
      (V_{OUT} - V_{cc})G_5 + I_{M_2} &= 0 \\
      V_{OUT}G_5 - V_{cc}G_5 + I_{M_2} &= 0 \\
      V_{OUT} = V_{cc} - \frac{I_{M_2}}{G_5} \tag{1}
    \end{align*}
    Applying KCL at $V_{X}$:
    \begin{align*}
      I_{M_0} + I_{G_1} - I_{M_3} &= 0 \\
      I_{M_0} + I_{G_1} = I_{M_3} \\
      I_{M_0} &\approx I_{M_3} \quad \text{[assuming } I_{G_1} \approx 0 \text{ by MOS physics]} \tag{2}
    \end{align*}
    Applying KCL at $V_{Y}$:
    \begin{align*}
      I_{M_1} - I_{M_2} &= 0 \\
      I_{M_1} = I_{M_2} \tag{3}
    \end{align*}
    Applying KCL at $V_{Z}$:
    \begin{align*}
      (V_{Z} - V_{cc})G_4 + I_{M_3} + I_{G_2} &= 0 \\
      V_{Z}G_4 - V_{cc}G_4 + I_{M_3} + I_{G_2} &= 0 \\
      V_{Z}G_4 &= V_{cc}G_4 - (I_{M_3} + I_{G_2}) \\
      V_{Z} &= V_{cc} - \frac{(I_{M_3} + I_{G_2})}{G_4} \tag{4}
    \end{align*}
    Using the transistor equations for $I_{M_0}$, $I_{M_1}$, $I_{M_2}$, and $I_{M_3}$:
    \begin{align*}
      I_{M_0} &= k_{n_0} (V_{GS_0} - V_{th})^2 (1 + K_{CLM} V_{DS_0}) \\
              &= k_{n_0} (V_{X} - V_{ee} - V_{th})^2 (1 + K_{CLM} (V_{X} - V_{ee})) \\
              &= k_{n_0} (V_{X} - V_{ee} - V_{th})^2 \quad \text{[assuming } K_{CLM} \approx 0 \text{]} \tag{5} \\
      I_{M_1} &= k_{n_1} (V_{GS_1} - V_{th})^2 (1 + K_{CLM} V_{DS_1}) \\
              &= k_{n_1} (V_{X} - V_{ee} - V_{th})^2 (1 + K_{CLM} (V_{Y} - V_{ee})) \\
              &= k_{n_1} (V_{X} - V_{ee} - V_{th})^2 \quad \text{[assuming } K_{CLM} \approx 0 \text{]} \tag{6} \\
      I_{M_2} &= k_{n_2} (V_{GS_2} - V_{th})^2 (1 + K_{CLM} V_{DS_2}) \\
              &= k_{n_2} (V_{Z} - V_{Y} - V_{th})^2 (1 + K_{CLM} (V_{OUT} - V_{Y})) \\
              &= k_{n_2} (V_{Z} - V_{Y} - V_{th})^2 \quad \text{[assuming } K_{CLM} \approx 0 \text{]} \tag{7} \\
      I_{M_3} &= k_{n_3} (V_{GS_3} - V_{th})^2 (1 + K_{CLM} V_{DS_3}) \\
              &= k_{n_3} (V_{Z} - V_{X} - V_{th})^2 (1 + K_{CLM} (V_{Z} - V_{X})) \\
              &= k_{n_3} (V_{Z} - V_{X} - V_{th})^2 \quad \text{[assuming } K_{CLM} \approx 0 \text{]} \tag{8}
    \end{align*}
    It can be seen that,
    \begin{align*}
      I_{IN} = I_{M_0} &= I_{M_3} \quad \text{[from (2)]} \\
      I_{OUT} = I_{M_1} &= I_{M_2} \quad \text{[from (3)]} \\
      I_{M_0} &= I_{M_1} \quad \text{[from (5) and (6) and assuming } k_{n_0} = k_{n_1} \text{]} \\
      \therefore, I_{IN} &= I_{OUT}
    \end{align*}
    Thus,
    \begin{align*}
      \frac{I_{OUT}}{I_{IN}} &= \frac{I_{M_1}}{I_{M_0}} = \frac{k_{n_1} (V_{X} - V_{ee} - V_{th})^2}{k_{n_0} (V_{X} - V_{ee} - V_{th})^2} \hspace{1cm} \text{(From (5), (6): same } V_{GS}\text{)} \\
      &= \frac{k_{n_1}}{k_{n_0}} = \frac{\frac{K_P}{2}\frac{W_1}{L_1}}{\frac{K_P}{2}\frac{W_0}{L_0}} = \frac{\frac{W_1}{L_1}}{\frac{W_0}{L_0}} = \frac{W_1 L_0}{L_1 W_0}
    \end{align*}
    \(M_2\) carries \(I_{M_1}\) and \(M_3\) carries \(I_{M_0}\) (from (2), (3)); choosing
    \(\frac{W_2 L_3}{L_2 W_3} = \frac{W_1 L_0}{L_1 W_0}\) gives \(M_2\) and \(M_3\) the same
    \(V_{ov}\) (from (7), (8)), so \(V_Y = V_Z - V_{th} - V_{ov_2} = V_X\): \(M_1\) then sees
    \(V_{DS_1} = V_{DS_0}\) and the ratio holds even with \(K_{CLM} \neq 0\).
  \item \textbf{Small Signal, Frequency Domain, AC Input Analysis (Q4)}:\\
    Applying KCL at $V_{out}$:
    \begin{align*}
      V_{out}G_5 + (V_{out} - V_{y})g_{o_2} + (V_{z} - V_{y})g_{m_2} &= 0 \\
      V_{out}G_5 + V_{out}g_{o_2} - V_{y}g_{o_2} + V_{z}g_{m_2} - V_{y}g_{m_2} &= 0 \\
      V_{out}(G_5 + g_{o_2}) - V_{y}(g_{o_2} + g_{m_2}) + V_{z}g_{m_2} &= 0 \tag{9}
    \end{align*}
    Applying KCL at $V_{x}$:
    \begin{align*}
      V_{x}g_{o_0} + I_{G_1} + V_{x}g_{m_0} + (V_{x} - V_{z})g_{o_3} - (V_{z} - V_{x})g_{m_3} &= 0 \\
      V_{x}g_{o_0} + V_{x}g_{m_0} + (V_{x} - V_{z})g_{o_3} - (V_{z} - V_{x})g_{m_3} &= 0 \quad \text{[as } I_{G} \approx 0 \text{ by MOS physics]} \\
      V_{x}g_{o_0} + V_{x}g_{m_0} + V_{x}g_{o_3} - V_{z}g_{o_3} - V_{z}g_{m_3} + V_{x}g_{m_3} &= 0 \\
      V_{x}(g_{o_0} + g_{m_0} + g_{o_3} + g_{m_3}) - V_{z}(g_{o_3} + g_{m_3}) &= 0 \tag{10}
    \end{align*}
    Applying KCL at $V_{y}$:
    \begin{align*}
      V_{y}g_{o_1} + V_{x}g_{m_1} + (V_{y} - V_{out})g_{o_2} - (V_{z} - V_{y})g_{m_2} &= 0 \\
      V_{y}g_{o_1} + V_{x}g_{m_1} + V_{y}g_{o_2} - V_{out}g_{o_2} - V_{z}g_{m_2} + V_{y}g_{m_2} &= 0 \\
      V_{y}(g_{o_1} + g_{o_2} + g_{m_2}) - V_{out}g_{o_2} - V_{z}g_{m_2} + V_{x}g_{m_1} &= 0 \tag{11}
    \end{align*}
    Applying KCL at $V_{z}$:
    \begin{align*}
      V_{z}G_4 + (V_{z} - V_{x})g_{o_3} + (V_{z} - V_{x})g_{m_3} + I_{G_2} &= 0 \\
      V_{z}G_4 + (V_{z} - V_{x})g_{o_3} + (V_{z} - V_{x})g_{m_3} &= 0 \quad \text{[as } I_{G} \approx 0 \text{ by MOS physics]} \\
      V_{z}G_4 + V_{z}g_{o_3} - V_{x}g_{o_3} + V_{z}g_{m_3} - V_{x}g_{m_3} &= 0 \\
      V_{z}(G_4 + g_{o_3} + g_{m_3}) - V_{x}(g_{o_3} + g_{m_3}) &= 0 \tag{12}
    \end{align*}
    \paragraph{Voltage Transfer \(\frac{V_{out}}{V_z}\):} a source \(V_s\) forcing \(V_z = V_s\)
    supplies whatever current (12) needs, so (12) drops out and (9), (10),
    (11) are solved with \(V_z\) known. Let
    \begin{align*}
      S &= g_{m_0} + g_{o_0} + g_{m_3} + g_{o_3} \hspace{1cm} P = g_{m_2} + g_{o_1} + g_{o_2} \hspace{1cm} a = \frac{g_{m_3} + g_{o_3}}{S}
    \end{align*}
    From (10):
    \begin{align*}
      V_x S &= V_z (g_{o_3} + g_{m_3}) \\
      V_x &= a V_z \tag{13}
    \end{align*}
    From (11) with (13):
    \begin{align*}
      V_y P &= V_{out} g_{o_2} + V_z g_{m_2} - a V_z g_{m_1} \\
      V_y &= \frac{V_{out} g_{o_2} + V_z (g_{m_2} - a g_{m_1})}{P} \tag{14}
    \end{align*}
    Substituting (14) into (9) and multiplying by \(P\):
    \begin{align*}
      V_{out}(G_5 + g_{o_2}) P - (g_{o_2} + g_{m_2})\left(V_{out} g_{o_2} + V_z (g_{m_2} - a g_{m_1})\right) + V_z g_{m_2} P &= 0 \\
      V_{out}\left[(G_5 + g_{o_2}) P - g_{o_2}(g_{o_2} + g_{m_2})\right] &= V_z\left[(g_{o_2} + g_{m_2})(g_{m_2} - a g_{m_1}) - g_{m_2} P\right]
    \end{align*}
    The left bracket, with \(P - (g_{o_2} + g_{m_2}) = g_{o_1}\):
    \begin{align*}
      G_5 P + g_{o_2} P - g_{o_2}(g_{o_2} + g_{m_2}) &= G_5 P + g_{o_2}\left(P - g_{o_2} - g_{m_2}\right) = G_5 P + g_{o_1} g_{o_2}
    \end{align*}
    The right bracket:
    \begin{align*}
      (g_{o_2} + g_{m_2}) g_{m_2} - a g_{m_1}(g_{o_2} + g_{m_2}) - g_{m_2} P
        &= -g_{m_2}\left(P - g_{o_2} - g_{m_2}\right) - a g_{m_1}(g_{m_2} + g_{o_2}) \\
        &= -g_{m_2} g_{o_1} - a g_{m_1}(g_{m_2} + g_{o_2})
    \end{align*}
    so, multiplying numerator and denominator by \(S\) (\(aS = g_{m_3} + g_{o_3}\)):
    \begin{align*}
      \frac{V_{out}}{V_z} &= -\frac{g_{m_2} g_{o_1} + a g_{m_1}(g_{m_2} + g_{o_2})}{G_5 P + g_{o_1} g_{o_2}} \\
      \frac{V_{out}}{V_z} &= -\frac{g_{m_1}(g_{m_2} + g_{o_2})(g_{m_3} + g_{o_3}) + g_{m_2} g_{o_1}(g_{m_0} + g_{o_0} + g_{m_3} + g_{o_3})}{(g_{m_0} + g_{o_0} + g_{m_3} + g_{o_3})\left(G_5 (g_{m_2} + g_{o_1} + g_{o_2}) + g_{o_1} g_{o_2}\right)}
    \end{align*}
    \paragraph{Current Transfer \(\frac{I_{out}}{I_{in}}\):} with the output shorted
    (\(V_{out} = 0\)), \(I_{in}\) is the current \(M_3\) delivers into \(V_x\) and \(I_{out}\) the
    current \(M_1\) draws from \(V_y\). KCL at \(V_x\) says all of \(I_{in}\) flows down
    \(M_0\), and KCL at \(V_y\) says all of \(I_{out}\) flows down \(M_1\):
    \begin{align*}
      I_{in} &= (g_{m_0} + g_{o_0}) V_x \hspace{1cm} I_{out} = g_{m_1} V_x + g_{o_1} V_y
    \end{align*}
    From (14) with \(V_{out} = 0\) and \(V_z = \frac{V_x}{a} = V_x\frac{S}{g_{m_3} + g_{o_3}}\) from (13):
    \begin{align*}
      V_y &= \frac{V_z g_{m_2} - V_x g_{m_1}}{P} = \frac{V_x}{P}\left(\frac{g_{m_2} S}{g_{m_3} + g_{o_3}} - g_{m_1}\right) \\
      I_{out} &= g_{m_1} V_x + \frac{g_{o_1} V_x}{P}\left(\frac{g_{m_2} S}{g_{m_3} + g_{o_3}} - g_{m_1}\right)
               = \frac{V_x}{P}\left(g_{m_1}(P - g_{o_1}) + \frac{g_{o_1} g_{m_2} S}{g_{m_3} + g_{o_3}}\right) \\
      I_{out} &= V_x\frac{g_{m_1}(g_{m_2} + g_{o_2})(g_{m_3} + g_{o_3}) + g_{m_2} g_{o_1} S}{P (g_{m_3} + g_{o_3})} \\
      \frac{I_{out}}{I_{in}} &= \frac{g_{m_1}(g_{m_2} + g_{o_2})(g_{m_3} + g_{o_3}) + g_{m_2} g_{o_1} S}{(g_{m_0} + g_{o_0})(g_{m_3} + g_{o_3})(g_{m_2} + g_{o_1} + g_{o_2})}
    \end{align*}
    With \(g_m \gg g_o\), keeping the largest term in each factor:
    \(\frac{I_{out}}{I_{in}} \approx \frac{g_{m_1} g_{m_2} g_{m_3}}{g_{m_0} g_{m_3} g_{m_2}} = \frac{g_{m_1}}{g_{m_0}}\), the
    Q5 ratio \(\frac{W_1 L_0}{L_1 W_0}\) for equal \(V_{ov}\).
\end{enumerate}


\subsubsection{Analysis}
\begin{enumerate}
  \item \textbf{Output Load (Max)}: \(M_2\) saturated needs
    \(V_{DS_2} \geq V_{ov_2}\). With \(V_{GS_2} = V_Z - V_Y = V_{th} + V_{ov_2}\), i.e.
    \(V_Y = V_Z - V_{th} - V_{ov_2}\):
    \begin{align*}
      V_{OUT,\min} &= V_Y + V_{ov_2} = V_Z - V_{th} - V_{ov_2} + V_{ov_2} = V_Z - V_{th} \\
      V_{cc} - I_{M_2} R_{L} &\geq V_{OUT,\min} \\
      R_{L, \max} &= \frac{V_{cc} - V_{OUT,\min}}{I_{M_2}}
    \end{align*}
  \item \textbf{Impedances}:
    \paragraph{At Input:} a test current \(I_{test}\) into \(V_z\) adds to the
    left side of (12):
      \begin{align*}
        V_{z}(G_4 + g_{o_3} + g_{m_3}) - V_{x}(g_{o_3} + g_{m_3}) &= I_{test} \tag{15}
      \end{align*}
      (10) still gives \(V_x = a V_z\) from (13), so
      \begin{align*}
        I_{test} &= V_z\left(G_4 + (g_{o_3} + g_{m_3}) - a(g_{o_3} + g_{m_3})\right) = V_z\left(G_4 + (g_{m_3} + g_{o_3})(1 - a)\right) \\
        1 - a &= \frac{S - (g_{m_3} + g_{o_3})}{S} = \frac{g_{m_0} + g_{o_0}}{S} \\
        I_{test} &= V_z\left(G_4 + \frac{(g_{m_3} + g_{o_3})(g_{m_0} + g_{o_0})}{S}\right) = V_z\frac{G_4 S + (g_{m_3} + g_{o_3})(g_{m_0} + g_{o_0})}{S} \\
        Z_{in} &= \frac{V_{z}}{I_{test}} = \frac{g_{m_0} + g_{m_3} + g_{o_0} + g_{o_3}}{G_{4}(g_{m_0} + g_{m_3} + g_{o_0} + g_{o_3}) + (g_{m_0} + g_{o_0})(g_{m_3} + g_{o_3})}
      \end{align*}
      i.e.\ \(R_4 \parallel \left(\frac{1}{g_{m_0} + g_{o_0}} + \frac{1}{g_{m_3} + g_{o_3}}\right)\): two diodes in series
      shunted by \(R_4\).
    \paragraph{At Output:} a test current \(I_{test}\) into \(V_{out}\) adds to
      the left side of (9):
      \begin{align*}
        V_{out}(G_5 + g_{o_2}) - V_{y}(g_{o_2} + g_{m_2}) + V_{z}g_{m_2} &= I_{test} \tag{16}
      \end{align*}
      (10) and (12) are now both homogeneous in \(V_x, V_z\); from (10)
      \(V_x = aV_z\), and (12) becomes
      \begin{align*}
        V_z\left(G_4 + (g_{m_3} + g_{o_3})(1 - a)\right) &= 0
      \end{align*}
      The bracket is positive, so \(V_z = 0\) and \(V_x = 0\). Then (11) gives
      \begin{align*}
        V_y P &= V_{out} g_{o_2} \hspace{1cm} \Rightarrow \hspace{1cm} V_y = \frac{g_{o_2}}{P} V_{out}
      \end{align*}
      and (16):
      \begin{align*}
        I_{test} &= V_{out}(G_5 + g_{o_2}) - \frac{g_{o_2}(g_{o_2} + g_{m_2})}{P} V_{out}
                  = V_{out}\frac{G_5 P + g_{o_2}(P - g_{o_2} - g_{m_2})}{P} = V_{out}\frac{G_5 P + g_{o_1} g_{o_2}}{P} \\
        Z_{out} &= \frac{V_{out}}{I_{test}} = \frac{g_{m_2} + g_{o_1} + g_{o_2}}{G_{5}(g_{m_2} + g_{o_1} + g_{o_2}) + g_{o_1} g_{o_2}}
      \end{align*}
\end{enumerate}


\subsubsection{Design}
Large Signal, Time Domain, DC Input Analysis (Q5) can be used to set the bias
currents and voltages in this circuit.
\begin{enumerate}
  \item Choose appropriate transistor dimensions
    ($W$ and $L$) for transistors $M_0$ and $M_1$
    to achieve the desired current gain.
    \begin{align*}
      \frac{I_{OUT}}{I_{IN}} &= \frac{W_1 L_0}{L_1 W_0} = \frac{W_2 L_3}{L_2 W_3}
    \end{align*}
  \item Choose a suitable overdrive voltage value
    ($V_{ov}$) for computing the bias voltages at
    nodes $V_{X}$, $V_{Y}$, and $V_{Z}$.
    \begin{align*}
      V_{ov_0} &= V_{GS_0} - V_{th} \\
      V_{ov_0} &= V_{X} - V_{ee} - V_{th} \\
      V_{X} &= V_{ov_0} + V_{ee} + V_{th} \\
      \\
      V_{ov_3} &= V_{GS_3} - V_{th} \\
      V_{ov_3} &= V_{Z} - V_{X} - V_{th} \\
      V_{Z} &= V_{ov_3} + V_{X} + V_{th} \\
      \\
      V_{ov_2} &= V_{GS_2} - V_{th} \\
      V_{ov_2} &= V_{Z} - V_{Y} - V_{th} \\
      V_{Y} &= V_{Z} - V_{ov_2} - V_{th}
    \end{align*}
  \item Using the chosen $V_{ov}$, compute $W_0$
    to set the required bias current $I_{IN}$.
    \begin{align*}
      I_{IN} &\approx I_{M_0} = \frac{K_P}{2} \frac{W_0}{L_0} V_{ov_0}^2 \\
             &\approx I_{M_3} = \frac{K_P}{2} \frac{W_3}{L_3} V_{ov_3}^2 \\
      \therefore W_0 &= \frac{2 I_{IN} L_0}{K_P V_{ov_0}^2} \\
                 W_3 &= \frac{2 I_{IN} L_3}{K_P V_{ov_3}^2}
    \end{align*}
    $W_1$ can then be computed using the current
    gain relation from step (a). Similarly, $W_2$
    can be computed using the current gain relation
    from step (a).
  \item Compute the required value of $G_4$ (or
    $R_4$) to ensure that the voltage at node
    $V_{Z}$ is as calculated in step (b).
    \begin{align*}
      V_{Z} &= V_{cc} - \frac{(I_{M_3} + I_{G_2})}{G_4} \\
      V_{Z} &\approx V_{cc} - \frac{I_{M_3}}{G_4} \quad \text{[assuming } I_{G_2} \approx 0 \text{ by MOS physics]} \\
      V_{Z} - V_{cc} &\approx - \frac{I_{M_3}}{G_4} \\
      \frac{(I_{M_3})}{(V_{cc} - V_{Z})} &\approx G_4 \\
      \frac{(I_{IN})}{(V_{cc} - V_{Z})} &\approx G_4 \quad \text{[} I_{IN} = I_{M_3} = I_{M_0} \text{]} \\
      \therefore R_4 &\approx \frac{V_{cc} - V_{Z}}{I_{IN}}
    \end{align*}
\end{enumerate}

\subsubsection{Question}
Design a cascode current mirror with a current
gain of 5. The input current is $I_{IN} = 10 \mu$A.
Assume the following specifications:
\begin{enumerate}
  \item $V_{cc} = +5$ V, $V_{ee} = 0$ V
  \item $V_{th} = 0.7$ V
  \item $K_P = 100 \mu A/V^2$
  \item $K_{CLM} = 0.02$ V$^{-1}$
  \item $L_{\min} = 1 \mu m$
  \item $V_{L, \max} = V_{cc}$
  \item $V_{ov} = 0.2$ V for all transistors
\end{enumerate}
Find the following:
\begin{enumerate}
  \item Transistor dimensions ($W$ and $L$) for all
    transistors.
  \item Values of resistor $R_4$ required to set
    the bias current.
  \item Output voltage range and maximum output
    load that can be connected at the output.
  \item Small signal output resistance of the
    current mirror. Assume a load resistance
    $R_L = 50 k\Omega$ is connected at the
    output.
\end{enumerate}
\paragraph{Solution:}
\begin{enumerate}
  \item Transistor dimensions:
    \begin{align*}
      \frac{I_{OUT}}{I_{IN}} &= \frac{W_1 L_0}{L_1 W_0} = \frac{W_2 L_3}{L_2 W_3} = 5 \\
      \\
      V_{X} &= V_{ov_0} + V_{ee} + V_{th} = 0.2 + 0 + 0.7 = 0.9 \text{ V} \\
      V_{Z} &= V_{ov_3} + V_{X} + V_{th} = 0.2 + 0.9 + 0.7 = 1.8 \text{ V} \\
      V_{Y} &= V_{Z} - V_{ov_2} - V_{th} = 1.8 - 0.2 - 0.7 = 0.9 \text{ V} \\
      \\
      W_0 &= \frac{2 I_{IN} L_0}{K_P V_{ov_0}^2} = \frac{2 \times 10 \times 10^{-6} \times 1 \times 10^{-6}}{100 \times 10^{-6} \times (0.2)^2} = 5 \mu m \\
      W_3 &= \frac{2 I_{IN} L_3}{K_P V_{ov_3}^2} = \frac{2 \times 10 \times 10^{-6} \times 1 \times 10^{-6}}{100 \times 10^{-6} \times (0.2)^2} = 5 \mu m \\
      W_1 &= 5 \times \frac{W_0}{L_0} \times L_1 = 5 \times \frac{5}{1} \times 1 = 25 \mu m \\
      W_2 &= 5 \times \frac{W_3}{L_3} \times L_2 = 5 \times \frac{5}{1} \times 1 = 25 \mu m
    \end{align*}
  \item Value of resistor $R_4$:
    \begin{align*}
      R_4 &\approx \frac{V_{cc} - V_{Z}}
                      {I_{IN}} = \frac{5 - 1.8}{10 \times 10^{-6}} = 320 k\Omega
    \end{align*}
  \item Output voltage range and maximum output load
    (\(I_{M_2} = I_{OUT} = 5 \cdot 10\,\mu A = 50\,\mu A\)):
    \begin{align*}
      V_{OUT,\min} &\approx V_{Y} + V_{ov_2} = 0.9 + 0.2 = 1.1 \text{ V} \\
      V_{OUT,\max} &\approx V_{cc} = 5 \text{ V} \\
      R_{L, \max} &= \frac{V_{cc} - V_{OUT,\min}}{I_{M_2}} = \frac{5 - 1.1}{50 \times 10^{-6}} = 78 k\Omega
    \end{align*}
    With \(R_L = 50\,k\Omega\): \(V_{OUT} = V_{cc} - I_{M_2} R_L = 5 - 50 \times 10^{-6} \times 50 \times 10^{3} = 5 - 2.5 = 2.5\,V \geq 1.1\,V\),
    so the \(50\,k\Omega\) load keeps \(M_2\) saturated.
  \item Small signal output resistance:
    \begin{align*}
      g_{m_2} &= \frac{2 I_{M_2}}{V_{ov_2}} = \frac{2 \times 50 \times 10^{-6}}{0.2} = 5 \times 10^{-4} \text{ S} \\
      g_{o_1} &= I_{M_1} K_{CLM} = 50 \times 10^{-6} \times 0.02 = 1 \times 10^{-6} \text{ S} \\
      g_{o_2} &= I_{M_2} K_{CLM} = 50 \times 10^{-6} \times 0.02 = 1 \times 10^{-6} \text{ S} \\
      G_5 &= \frac{1}{R_L} = \frac{1}{50 k\Omega} = 2 \times 10^{-5} \text{ S} \\
      \\
      g_{m_2} + g_{o_1} + g_{o_2} &= 5 \times 10^{-4} + 1 \times 10^{-6} + 1 \times 10^{-6} = 5.02 \times 10^{-4} \text{ S} \\
      Z_{out} &= \frac{g_{m_2} + g_{o_1} + g_{o_2}}{G_{5}(g_{m_2} + g_{o_1} + g_{o_2}) + g_{o_1} g_{o_2}} \\
              &= \frac{5.02 \times 10^{-4}}{(2 \times 10^{-5})(5.02 \times 10^{-4}) + (10^{-6})(10^{-6})} \\
              &= \frac{5.02 \times 10^{-4}}{1.004 \times 10^{-8} + 10^{-12}} = \frac{5.02 \times 10^{-4}}{1.0041 \times 10^{-8}} \approx 49.995 k\Omega
    \end{align*}
    The load dominates. The mirror alone (\(G_5 \to 0\)):
    \begin{align*}
      Z_{out,mirror} &= \frac{g_{m_2} + g_{o_1} + g_{o_2}}{g_{o_1} g_{o_2}} = \frac{5.02 \times 10^{-4}}{10^{-12}} = 502 M\Omega \\
      g_{m_2} r_{o_1} r_{o_2} &= \frac{g_{m_2}}{g_{o_1} g_{o_2}} = \frac{5 \times 10^{-4}}{10^{-12}} = 500 M\Omega \hspace{1cm} \text{(the \(g_m \gg g_o\) estimate)}
    \end{align*}
\end{enumerate}

\subsubsection{Code}
\begin{lstlisting}[language=Octave]
  clear; clc;
  pkg load symbolic;

  syms Vx Vy Vz Vout Itest
  syms G5 G4 go0 gm0 go1 gm1 go2 gm2 go3 gm3

  % small signal kcl, equations (9) to (12)
  eq9  = Vout*(G5 + go2) - Vy*(go2 + gm2) + Vz*gm2 == 0;  % KCL at Vout
  eq10 = Vx*(go0 + gm0 + go3 + gm3) - Vz*(go3 + gm3) == 0; % KCL at Vx
  eq11 = Vy*(go1 + go2 + gm2) - Vout*go2 - Vz*gm2 + Vx*gm1 == 0; % KCL at Vy
  eq12 = Vz*(G4 + go3 + gm3) - Vx*(go3 + gm3) == 0;        % KCL at Vz

  % voltage transfer: Vz forced by a source, so (12) drops out
  sol = solve([eq9, eq10, eq11], [Vout, Vx, Vy]);
  printf('=== Vout/Vz ===\n'); pretty(factor(simplify(sol.Vout/Vz)))

  % current transfer: output shorted (Vout = 0), Vz forced
  sol = solve([eq10, subs(eq11, Vout, 0)], [Vx, Vy]);
  Iin  = (gm0 + go0)*sol.Vx;            % all of Iin flows down M0
  Iout = gm1*sol.Vx + go1*sol.Vy;       % all of Iout flows down M1
  printf('=== Iout/Iin ===\n'); pretty(factor(simplify(Iout/Iin)))

  % input impedance: Itest injected into Vz, equation (15)
  eq15 = Vz*(G4 + go3 + gm3) - Vx*(go3 + gm3) == Itest;
  sol = solve([eq9, eq10, eq11, eq15], [Vout, Vx, Vy, Vz]);
  printf('=== Zin ===\n'); pretty(factor(simplify(sol.Vz/Itest)))

  % output impedance: Itest injected into Vout, equation (16)
  eq16 = Vout*(G5 + go2) - Vy*(go2 + gm2) + Vz*gm2 == Itest;
  sol = solve([eq16, eq10, eq11, eq12], [Vout, Vx, Vy, Vz]);
  printf('=== Zout ===\n'); pretty(factor(simplify(sol.Vout/Itest)))
\end{lstlisting}

